\documentclass[8pt,a4paper]{article}

% ------------------------------------------------------------------
% Packages
% ------------------------------------------------------------------
\usepackage[utf8]{inputenc}
\usepackage[T1]{fontenc}
\usepackage{lmodern}
\usepackage[english]{babel}
\usepackage{amsmath, amssymb, amsthm, mathtools}
\usepackage[margin=2.5cm]{geometry}
\usepackage{hyperref}
\usepackage{graphicx}
\usepackage{xcolor}
\usepackage{listings}
\usepackage{booktabs}
\usepackage{enumitem}
\usepackage{siunitx}
\usepackage{bm}

\hypersetup{
  colorlinks=true,
  linkcolor=blue!60!black,
  citecolor=blue!60!black,
  urlcolor=blue!60!black
}

% ------------------------------------------------------------------
% Notation shortcuts
% ------------------------------------------------------------------
\newcommand{\R}{\mathbb{R}}
\newcommand{\bvec}[1]{\bm{#1}}
\newcommand{\xvec}{\bvec{x}}
\newcommand{\uvec}{\bvec{u}}
\newcommand{\fvec}{\bvec{f}}
\newcommand{\lamr}{\bvec{\lambda}_r}
\newcommand{\lamv}{\bvec{\lambda}_v}
\newcommand{\Omegax}{\Omega_x}
\newcommand{\Omegay}{\Omega_y}
\newcommand{\Omegaxx}{\Omega_{xx}}
\newcommand{\Omegayy}{\Omega_{yy}}
\newcommand{\Omegaxy}{\Omega_{xy}}
\newcommand{\Amat}{\bvec{A}}
\newcommand{\Bmat}{\bvec{B}}
\newcommand{\Qmat}{\bvec{Q}}
\newcommand{\Vr}{\bvec{V}_r}
\newcommand{\Ddec}{\bvec{D}_{\mathrm{dec}}}
\newcommand{\Fmat}{\bvec{F}}
\newcommand{\Hmat}{\bvec{H}}
\newcommand{\TT}{^{\mathsf{T}}}

\lstset{
  basicstyle=\ttfamily\footnotesize,
  frame=single,
  numbers=left,
  numberstyle=\tiny\color{gray},
  breaklines=true,
  tabsize=2
}

\theoremstyle{definition}
\newtheorem{definition}{Definition}[section]

\theoremstyle{remark}
\newtheorem*{remark}{Remark}

% ------------------------------------------------------------------
\title{Quantum Annealing for Earth--Moon Low-Thrust Trajectory Optimization}
\author{Gino Luciano Moretta \\
  \small Sapienza Universit\`a di Roma, supervisor Prof.\ A.\ Paolozzi, E. Ortore}
\date{Last updated: April 2026}

\begin{document}
\maketitle
\tableofcontents
\bigskip

% ==================================================================
\section{Project overview}
% ==================================================================

\paragraph{Research goal.}
Extend the pseudospectral quantum-annealing transcription framework of
De~Grossi, Carbone et~al.\ (\emph{Astrodynamics} 9:195--215,
2025)~\cite{Carbone2025} from Earth--Mars interplanetary low-thrust
transfers to \textbf{Earth--Moon} low-thrust transfers in the planar
circular restricted three-body problem (CR3BP), and evaluate the
resulting problem on D-Wave's \textbf{Advantage 2} system (and its
nonlinear hybrid solver) as part of a six-month research project in
aerospace engineering at Sapienza.

\paragraph{Companion paper.}
The scientific account of the on/off scheduling formulation, the
iterative re-linearisation methodology, the classical
mixed-integer-linear-programming baseline, and the experimental
results on the cislunar injection-correction scenario is given in
the companion paper, \emph{Iterative Re-Linearisation for a
Naturally-Binary QUBO Formulation of Low-Thrust Cislunar Trajectory
Scheduling} (\texttt{paper/paper.tex}). This manual is the
technical reference for the implementation: it covers the full
codebase end-to-end, including the earlier continuous-coefficient
QUBO that the project pivoted away from, and is intended for
developers and reviewers reading the source.

\paragraph{Why the Earth--Moon case is harder.}
Carbone et al.\ worked in an \emph{inertial two-body}
(Keplerian~$+$~small~perturbations) setting where the dynamics are
essentially linear around the reference heliocentric arc. The
Earth--Moon CR3BP in the rotating (\textbf{synodic}) frame is
\emph{genuinely nonlinear}: the gradient of the pseudo-potential
contains $1/r^2$ terms that blow up near each primary. A single-shot
least-squares $\rightarrow$ QUBO pipeline cannot accommodate this
directly, so we layer a \textbf{sequential linearization outer loop}
on top of the QUBO solver, at each iteration the nonlinear
dynamics are linearized about the current best trajectory, a new
linear system is built, and the QUBO solver is called again.

\paragraph{Two transcription philosophies.}
There are two principled ways to convert an optimal control problem
into a linear residual system that can be fed into a QUBO:

\begin{itemize}[leftmargin=1.5em, itemsep=2pt]
  \item \textbf{Direct}: discretize state and control at a set of
    collocation nodes (e.g.\ Legendre--Gauss--Lobatto), and write the
    dynamics constraints as algebraic residuals. Carbone et al.\ use
    this approach with LGL nodes and a Lagrange-polynomial
    differentiation matrix.
  \item \textbf{Indirect}: apply Pontryagin's Minimum Principle
    analytically to eliminate the control (for an energy-minimizing
    cost, $\uvec^\star = -\lamv$), obtain the two-point boundary value
    problem for state \emph{and} costate, and then discretize by a
    trial-function expansion. Gino's existing v2 rendezvous code uses
    this path with \textbf{Theory of Functional Connections (TFC)}
    switching functions plus an \textbf{Extreme Learning Machine
    (ELM)} basis for the free part of the solution.
\end{itemize}

This project will implement \textbf{both} and compare them on the
same Earth--Moon benchmark. The indirect TFC+ELM pipeline is
implemented first because it reuses almost all of Gino's validated
rendezvous machinery; the direct LGL pipeline is implemented second
as a baseline.


% ==================================================================
\section{Units and the planar CR3BP}
% ==================================================================

\subsection{Nondimensionalization}

We work in \textbf{canonical CR3BP units}. Let $m_1$ (Earth) and
$m_2$ (Moon) be the two primaries, $a$ the semi-major axis of their
mutual orbit, and $n$ the mean motion. Define:

\begin{itemize}[leftmargin=1.5em, itemsep=1pt]
  \item \textbf{Length unit} $L = a \approx 384\,400~\si{km}$.
  \item \textbf{Time unit} $T = 1/n$; one nondimensional time unit
    corresponds to roughly $4.343~\si{days} \approx 3.757 \times 10^5~\si{s}$.
  \item \textbf{Mass unit} $m_1 + m_2$, so the mass parameter is
    \[
      \mu \;=\; \frac{m_2}{m_1 + m_2}
          \;\approx\; 0.012\,150\,583\,925\,359
          \quad (\text{Earth--Moon}).
    \]
\end{itemize}

With this scaling, the gravitational constant times total mass is
unity, and the angular velocity of the rotating (synodic) frame is
$\omega = 1$. The primaries become \emph{stationary} on the $x$-axis:

\[
  \bvec{r}_{\oplus}^{\mathrm{syn}} = (-\mu,\, 0),
  \qquad
  \bvec{r}_{\text{Moon}}^{\mathrm{syn}} = (1 - \mu,\, 0).
\]

\begin{remark}
This is why the CR3BP is such a clean playground for mission design:
in the synodic frame, geometry is frozen, and the dynamics become a
time-invariant vector field on a 4D phase space.
\end{remark}

\subsection{Equations of motion}

Let $\xvec = (x, y)$ be the synodic position and $\uvec = (u_x, u_y)$
the thrust-acceleration control (low-thrust is modelled as a
continuous acceleration, not an impulsive $\Delta v$). The planar
CR3BP with control reads

\begin{equation}
\boxed{
  \begin{aligned}
    \ddot{x} &= \phantom{-}2\dot{y} + \Omegax(x, y) + u_x, \\
    \ddot{y} &=           -2\dot{x} + \Omegay(x, y) + u_y,
  \end{aligned}
}
\label{eq:cr3bp}
\end{equation}

where the pseudo-potential $\Omega$ and its gradient are

\begin{align}
  \Omega(x, y) &= \tfrac12 (x^2 + y^2) + \frac{1-\mu}{r_1}
     + \frac{\mu}{r_2} + \tfrac12 \mu(1-\mu), \\[4pt]
  \Omegax(x,y) &= x - \frac{(1-\mu)(x+\mu)}{r_1^3}
     - \frac{\mu(x - 1 + \mu)}{r_2^3}, \\
  \Omegay(x,y) &= y - \frac{(1-\mu)\,y}{r_1^3} - \frac{\mu\,y}{r_2^3},
\end{align}

with distances from the two primaries

\[
  r_1 = \sqrt{(x+\mu)^2 + y^2}, \qquad
  r_2 = \sqrt{(x - 1 + \mu)^2 + y^2}.
\]

The two $2\dot{y}$ and $-2\dot{x}$ terms are the \textbf{Coriolis
coupling} from the rotating frame; the centrifugal contribution has
been absorbed into the $\tfrac12(x^2+y^2)$ part of $\Omega$.

In first-order form with state $\xvec = (x, y, v_x, v_y)$:

\begin{equation}
  \dot{\xvec} \;=\; \fvec(\xvec) + \Bmat \uvec,
  \qquad
  \Bmat \;=\;
  \begin{pmatrix}
    0 & 0 \\ 0 & 0 \\ 1 & 0 \\ 0 & 1
  \end{pmatrix},
  \qquad
  \fvec(\xvec) \;=\;
  \begin{pmatrix}
    v_x \\ v_y \\
    2 v_y + \Omegax \\
    -2 v_x + \Omegay
  \end{pmatrix}.
\label{eq:cr3bp-1st-order}
\end{equation}

This is the form consumed by \texttt{PlanarCR3BP.rhs} in the code.

\subsection{The Jacobi integral}

Along any unforced ($\uvec = \bm 0$) trajectory, the CR3BP conserves
the \textbf{Jacobi integral}

\begin{equation}
  C_J(\xvec) \;=\; 2\,\Omega(x, y) \;-\; v_x^2 \;-\; v_y^2.
  \label{eq:jacobi}
\end{equation}

This is the CR3BP's analogue of mechanical energy (in the synodic
frame) and is the single most useful sanity check for any integrator
or transcription: a pure-drift propagator that lets $C_J$ wander
immediately gives away a bug. The code verifies $C_J$ conservation to
$<10^{-8}$ over $10^3$ RK4 steps as a unit test.

\subsection{Lagrange points}

Equation~\eqref{eq:cr3bp} has \textbf{five} equilibria in the rotating
frame --- the five Lagrange points $L_1, L_2, L_3$ (collinear, on the
$x$-axis) and $L_4, L_5$ (triangular, on the equilateral triangles
with the primaries). For Earth--Moon:

\begin{center}
\begin{tabular}{lccc}
  \toprule
  Point & Approx.\ $x$ & Approx.\ $y$ & Stability (linear) \\
  \midrule
  $L_1$ & $0.8369$ & $0$ & saddle--center \\
  $L_2$ & $1.1557$ & $0$ & saddle--center \\
  $L_3$ & $-1.0051$ & $0$ & saddle--center \\
  $L_4$ & $0.5 - \mu$ & $+\sqrt{3}/2$ & stable ($\mu < 0.0385$) \\
  $L_5$ & $0.5 - \mu$ & $-\sqrt{3}/2$ & stable \\
  \bottomrule
\end{tabular}
\end{center}

The triangular points $L_4$ and $L_5$ sit at the apex of the
equilateral triangles $\oplus\text{--}\text{Moon}\text{--}L_{4/5}$,
i.e.\ at $r_1 = r_2 = 1$, and --- because of this and a minor bookkeeping
cancellation --- the Jacobi integral there is \emph{exactly} $C_J = 3$
independently of $\mu$.

\subsection{Linearization of~\eqref{eq:cr3bp}}
\label{sec:linearization}

The QUBO pipeline downstream is fundamentally a \emph{linear}
least-squares engine; it needs the dynamics to be linear in the
unknowns. For the CR3BP we therefore linearize~\eqref{eq:cr3bp} about
a \textbf{nominal trajectory} $\bar{\xvec}(t) = (\bar x(t), \bar
y(t), \bar v_x(t), \bar v_y(t))$ and solve for the
\emph{deviation} $\delta\xvec = \xvec - \bar{\xvec}$.

Taylor-expanding $\Omegax$ and $\Omegay$:

\begin{equation}
  \begin{aligned}
    \Omegax(x, y) &\approx \Omegax(\bar x, \bar y)
       + \Omegaxx(\bar x, \bar y)\,\delta x
       + \Omegaxy(\bar x, \bar y)\,\delta y, \\
    \Omegay(x, y) &\approx \Omegay(\bar x, \bar y)
       + \Omegaxy(\bar x, \bar y)\,\delta x
       + \Omegayy(\bar x, \bar y)\,\delta y,
  \end{aligned}
\end{equation}

where the Hessian components are

\begin{equation}
  \Omegaxx = 1 - \frac{1-\mu}{r_1^3} - \frac{\mu}{r_2^3}
       + \frac{3(1-\mu)(x+\mu)^2}{r_1^5}
       + \frac{3\mu(x - 1 + \mu)^2}{r_2^5},
\label{eq:omega-xx}
\end{equation}
\begin{equation}
  \Omegayy = 1 - \frac{1-\mu}{r_1^3} - \frac{\mu}{r_2^3}
       + \frac{3(1-\mu)\,y^2}{r_1^5} + \frac{3\mu\,y^2}{r_2^5},
\label{eq:omega-yy}
\end{equation}
\begin{equation}
  \Omegaxy = \frac{3(1-\mu)(x+\mu)\,y}{r_1^5}
       + \frac{3\mu(x - 1 + \mu)\,y}{r_2^5}.
\label{eq:omega-xy}
\end{equation}

These come from \texttt{PlanarCR3BP.jacobian\_state}, which is tested
against central finite differences across six random seeds.

Substituting the expansion into~\eqref{eq:cr3bp} and subtracting the
nominal equation yields a linear time-varying (LTV) system for
$\delta\xvec$:

\begin{equation}
  \boxed{
  \begin{aligned}
    \delta\ddot{x} &=  2\,\delta\dot{y}
       + \Omegaxx(t)\,\delta x + \Omegaxy(t)\,\delta y + \delta u_x, \\
    \delta\ddot{y} &= -2\,\delta\dot{x}
       + \Omegaxy(t)\,\delta x + \Omegayy(t)\,\delta y + \delta u_y,
  \end{aligned}
  }
  \label{eq:cr3bp-linearized}
\end{equation}

where the coefficients $\Omegaxx(t) \equiv \Omegaxx(\bar x(t),
\bar y(t))$ etc.\ are \emph{known} functions of time once the
nominal is fixed.

\begin{remark}
Equation~\eqref{eq:cr3bp-linearized} should look familiar: with the
nominal frozen at $L_4$ and Earth--Moon $\mu$, it reduces exactly to
the classical L4 libration linear analysis. With the nominal frozen
at the Earth, it reduces to Clohessy--Wiltshire-like behaviour (up to
a different reference frame). Both are physically meaningful
sanity-check regimes for the transcription.
\end{remark}


% ==================================================================
\section{Optimal control formulation}
% ==================================================================

\subsection{The problem}

We want to steer the spacecraft from an initial \textbf{low Earth orbit
(LEO)} state to a final \textbf{low lunar orbit (LLO)} state over a
fixed time of flight $T$:

\begin{equation}
  \min_{\uvec(\cdot)} \quad
    J \;=\; \tfrac12 \int_0^T \|\uvec(t)\|^2\, dt
\end{equation}
subject to the CR3BP dynamics~\eqref{eq:cr3bp-1st-order} and the
boundary conditions

\begin{equation}
  \bvec{r}(0) = \bvec{r}_0,\quad
  \bvec{v}(0) = \bvec{v}_0,\quad
  \bvec{r}(T) = \bvec{r}_f,\quad
  \bvec{v}(T) = \bvec{v}_f.
\label{eq:bcs}
\end{equation}

The $\tfrac12 \|\uvec\|^2$ cost is the standard \textbf{energy-optimal}
objective for low-thrust rendezvous. It is \emph{not} fuel-optimal
($\int \|\uvec\|\,dt$), but it is quadratic in the control so it stays
friendly to QUBO and least-squares pipelines. Fuel-optimal cost is an
explicit future extension.

\subsection{Hamiltonian and Pontryagin}

Introduce the costate $\bvec{\lambda}(t) = (\lamr, \lamv)^\top$,
partitioned so that $\lamr$ is conjugate to $\bvec{r}$ and $\lamv$
is conjugate to $\bvec{v}$. The control Hamiltonian is

\begin{equation}
  \mathcal{H}(\xvec, \uvec, \bvec{\lambda})
   \;=\;
   \tfrac12 \|\uvec\|^2
   + \lamr\TT \bvec{v}
   + \lamv\TT \Big( \bvec{g}(\bvec{r}) + \bvec{C} \bvec{v} + \uvec \Big),
\end{equation}

where $\bvec{g} = (\Omegax, \Omegay)$ is the pseudo-potential
gradient and

\begin{equation}
  \bvec{C} \;=\;
  \begin{pmatrix}
    0 & 2 \\ -2 & 0
  \end{pmatrix}
\end{equation}

is the Coriolis block (constant). Pontryagin's Minimum Principle
gives the optimality conditions:

\begin{align}
  \dot{\bvec{r}} &=\phantom{-} \partial \mathcal{H} / \partial \lamr
              \;=\; \bvec{v}, \\
  \dot{\bvec{v}} &=\phantom{-} \partial \mathcal{H} / \partial \lamv
              \;=\; \bvec{g}(\bvec{r}) + \bvec{C} \bvec{v} + \uvec, \\
  \dot{\lamr}    &=          - \partial \mathcal{H} / \partial \bvec{r}
              \;=\; - \Big(\partial \bvec{g}/\partial \bvec{r}\Big)\TT \lamv, \\
  \dot{\lamv}    &=          - \partial \mathcal{H} / \partial \bvec{v}
              \;=\; - \lamr - \bvec{C}\TT \lamv,
\end{align}
and the optimality equation
\begin{equation}
  0 \;=\; \partial \mathcal{H}/\partial \uvec
    \;=\; \uvec + \lamv
  \qquad\Longrightarrow\qquad
  \boxed{\,\uvec^\star \;=\; -\lamv.\,}
\end{equation}

\paragraph{Consequence.}
The control is \emph{eliminated analytically}. The two-point BVP for
the indirect method has unknowns $(\bvec{r}, \bvec{v}, \lamr, \lamv)$
--- $8$ scalar functions in the planar case --- with $8$ boundary
conditions (four from the BCs~\eqref{eq:bcs} on $\bvec{r}$, four on
$\bvec{v}$). The costates are \emph{unconstrained} at the boundaries:
they adjust automatically to whatever keeps the optimality conditions
consistent.

\subsection{Linearized costate equations}

When the state dynamics are linearized about a nominal as
in~\eqref{eq:cr3bp-linearized}, the costate dynamics linearize
accordingly because the Hessian of $\Omega$ is what appears in the
state Jacobian. Writing out all four components:

\begin{align}
  \dot\lambda_{r_x} &= -\big(\Omegaxx \lambda_{v_x}
                           + \Omegaxy \lambda_{v_y}\big), \\
  \dot\lambda_{r_y} &= -\big(\Omegaxy \lambda_{v_x}
                           + \Omegayy \lambda_{v_y}\big), \\
  \dot\lambda_{v_x} &= -\lambda_{r_x} + 2\,\lambda_{v_y}, \\
  \dot\lambda_{v_y} &= -\lambda_{r_y} - 2\,\lambda_{v_x},
\end{align}

where all $\Omega$ Hessian coefficients are evaluated at the nominal
trajectory (so they are \emph{known functions of time}). These are
the equations the transcription will discretize.


% ==================================================================
\section{Indirect TFC+ELM transcription}
\label{sec:transcription}
% ==================================================================

\subsection{Theory of Functional Connections in one picture}

The \textbf{Theory of Functional Connections} (TFC), due to
Mortari~\cite{Mortari2017}, gives a way to write down a general
functional form that \emph{automatically satisfies any set of linear
boundary conditions}. The basic idea:

\begin{quote}
\emph{Instead of enforcing boundary conditions as separate equations,
build them into the basis.} The basis becomes a family of
``constrained expressions'' indexed by free parameters $\xvec$,
and every choice of $\xvec$ lands on a function that already meets
the BCs by construction.
\end{quote}

For a second-order ODE on $[0, T]$ with four BCs on position and
velocity,

\begin{equation}
  x(0) = x_0,\quad x(T) = x_f,\quad \dot x(0) = v_0,\quad \dot x(T) = v_f,
\end{equation}

the constrained expression takes the form

\begin{equation}
  x(t) \;=\; C(t) \;+\; F(t, \xvec),
\end{equation}

where $C(t)$ is a particular solution hitting the BCs and $F(t,
\xvec)$ is the free part \emph{annihilated by all four BC
projections}. TFC provides the following standard choice in terms
of switching functions $\omega_1, \omega_2, \omega_3, \omega_4$:

\begin{align}
  C(t) &= \omega_1(t)\, x_0 + \omega_2(t)\, x_f
        + \omega_3(t)\, v_0 + \omega_4(t)\, v_f, \\
  F(t, \xvec) &= g(t, \xvec)
        - \omega_1(t)\, g(0, \xvec) - \omega_2(t)\, g(T, \xvec)
        - \omega_3(t)\, \dot g(0, \xvec) - \omega_4(t)\, \dot g(T, \xvec),
\end{align}

where $g(\cdot, \xvec)$ is \emph{any} twice-differentiable function
of $t$ parameterised by unknowns $\xvec$. The four $\omega_i$ are
cubic Hermite basis polynomials on $[0, T]$ with the following
``Kronecker-like'' property:

\begin{equation}
  \begin{array}{c|cccc}
     & \omega_1 & \omega_2 & \omega_3 & \omega_4 \\
    \hline
     \text{value at } 0           & 1 & 0 & 0 & 0 \\
     \text{value at } T           & 0 & 1 & 0 & 0 \\
     \text{derivative at } 0      & 0 & 0 & 1 & 0 \\
     \text{derivative at } T      & 0 & 0 & 0 & 1
  \end{array}
\end{equation}

so that $C(0) = x_0$, $C(T) = x_f$, $\dot C(0) = v_0$, $\dot C(T) =
v_f$, and $F$ contributes $0$ to all four BC projections for any
$\xvec$. It follows that the full expression satisfies all four BCs
\emph{independently} of $\xvec$ --- the optimizer can search freely
in the space of $\xvec$ without ever violating a BC.

\paragraph{Explicit formulas.}
On the normalized interval $z \in [-1, 1]$ with $dz = 2$, the
implementation uses

\begin{align}
  \omega_1(z) &=  1 + 2\tau^3/8 - 3\tau^2/4, \\
  \omega_2(z) &= -2\tau^3/8 + 3\tau^2/4, \\
  \omega_3(z) &=  \tau + \tau^3/4 - \tau^2, \\
  \omega_4(z) &=  \tau^3/4 - \tau^2/2,
\end{align}

with $\tau = z - z_0 = z + 1$. These are the same Hermite basis
polynomials used in Gino's rendezvous code; a time-rescaling factor
$c = (z_f - z_0)/(t_f - t_0) = 2/T$ converts derivatives with
respect to $z$ back into derivatives with respect to physical time.

\subsection{Extreme Learning Machine basis}

The ``free'' function $g(t, \xvec)$ that lives inside $F$ can be
\emph{any} twice-differentiable parameterization. TFC experiments
typically use a \textbf{single-hidden-layer neural network with
random, frozen input weights} --- an \textbf{Extreme Learning Machine}
(ELM)~\cite{Huang2006}:

\begin{equation}
  g(z, \xvec) \;=\; \sum_{j=1}^{L} \xi_j \, \sigma\big( w_j z + b_j \big),
\end{equation}

where $\sigma$ is a smooth activation (e.g.\ $\tanh$, logistic, or
Gaussian), the weights $w_j$ and biases $b_j$ are drawn \emph{once}
from a fixed distribution (e.g.\ uniform on $[-3, 3]$) and \emph{never
updated}, and the output weights $\xi_j$ are the only free parameters.

An ELM is a universal approximator: for any sufficiently smooth
target function $g^\star$ and any error tolerance $\varepsilon$, there
exists an ELM with enough neurons $L$ that gets within $\varepsilon$
of $g^\star$ by an appropriate choice of $\xvec$. Crucially, because
the output layer is linear in $\xvec$, the LSQ problem ``fit
$g(\cdot, \xvec)$ to residual data'' is linear in $\xvec$ --- exactly
what we need for the downstream QUBO engine.

\subsubsection*{Why not spline or Chebyshev?}

Two reasons ELM is preferred here:

\begin{enumerate}[leftmargin=1.5em]
  \item \textbf{Random-feature universality.} There is no rigid grid
    and no dependence on collocation structure, you simply draw
    more neurons to get more capacity.
  \item \textbf{Well-conditioning of the transcription matrix.}
    Gaussian and tanh ELM features are effectively random Fourier-like
    features: the linear system assembled from them typically has a
    smoothly decaying singular spectrum that plays well with the
    truncated-SVD step of the QUBO pipeline.
\end{enumerate}

\subsubsection*{Derivatives of the ELM basis}

At each training point $z_i$ and each neuron $j$, we need three
quantities: $\sigma$, $\sigma'$, $\sigma''$. For $\sigma = \tanh$:

\begin{align}
  \sigma(w z + b) &= \tanh(w z + b), \\
  \sigma'(w z + b) &= w\big(1 - \tanh^2(w z + b)\big), \\
  \sigma''(w z + b) &= -2 w^2 \tanh(w z + b)\big(1 - \tanh^2(w z + b)\big).
\end{align}

Stacking these over $n$ training points and $L$ neurons gives three
matrices $\Hmat, \Hmat', \Hmat'' \in \R^{n \times L}$.

\subsection{Constrained expression assembly}

With the switching functions and the ELM basis in hand, build the
following per-coordinate constrained matrices on the $n$ training
points:

\begin{align}
  \Fmat     &= \Hmat
              - \bvec{\omega}_1 \Hmat(z_0)^{\phantom\top}
              - \bvec{\omega}_2 \Hmat(z_f)
              - \bvec{\omega}_3 \Hmat'(z_0)
              - \bvec{\omega}_4 \Hmat'(z_f), \\
  \Fmat'    &= c\,\Big[\,\Hmat'
              - \bvec{\omega}_1' \Hmat(z_0)
              - \bvec{\omega}_2' \Hmat(z_f)
              - \bvec{\omega}_3' \Hmat'(z_0)
              - \bvec{\omega}_4' \Hmat'(z_f)\Big], \\
  \Fmat''   &= c^2\,\Big[\,\Hmat''
              - \bvec{\omega}_1'' \Hmat(z_0) - \dots\Big].
\end{align}

The factor $c = 2/T$ is the chain-rule conversion between $z$-time
and physical time, and appears once per derivative order.

Note the crucial property: $\Fmat$ is zero in the first and last
rows by construction (the switching-function subtraction exactly
cancels the value at the endpoints), and likewise $\Fmat'$ is zero
at the endpoints. This encodes the BC annihilation.

The BC-particular part is assembled similarly:

\begin{equation}
  \bvec{C}(z) \;=\; \bvec{\omega}_1 r_0 + \bvec{\omega}_2 r_f
      + \tfrac{1}{c}\big[\bvec{\omega}_3 v_0 + \bvec{\omega}_4 v_f\big],
\end{equation}

with the $1/c$ for velocity BCs because velocities in physical time
are the $z$-derivative times $c$. Evaluating $\bvec{C}$ and its
derivatives at each training point gives three vectors $\bvec{C},
\bvec{C}', \bvec{C}'' \in \R^{n}$.

\subsection{The linear residual system}
\label{sec:linear-system}

Gather all the unknown ELM weights into a single tall vector:

\begin{equation}
  \bm{\xi} \;=\; \big(
    \bm{\xi}_{x},\;
    \bm{\xi}_{y},\;
    \bm{\xi}_{\lambda_{r_x}},\;
    \bm{\xi}_{\lambda_{r_y}},\;
    \bm{\xi}_{\lambda_{v_x}},\;
    \bm{\xi}_{\lambda_{v_y}}
  \big) \in \R^{6L},
\end{equation}

so there are six blocks of $L$ weights --- two for the position
states (which use $\Fmat$-based constrained expressions with BCs)
and four for the costates (which use $\Hmat$-based
expressions \emph{without} BC enforcement, since the costates have
no fixed boundary values).

At each of the $n$ training points, the six ODEs yield $6n$ linear
residual equations. Written as a block matrix:

\begin{equation}
  \underbrace{\begin{pmatrix}
    \Fmat'' - \Omegaxx \Fmat & -\Omegaxy \Fmat  & \bm 0   & \bm 0 & \Hmat  & \bm 0  \\
    -\Omegaxy \Fmat & \Fmat'' - \Omegayy \Fmat  & \bm 0   & \bm 0 & \bm 0  & \Hmat  \\
    \bm 0 & \bm 0 & \Hmat' & \bm 0 & \Omegaxx \Hmat & \Omegaxy \Hmat \\
    \bm 0 & \bm 0 & \bm 0 & \Hmat' & \Omegaxy \Hmat & \Omegayy \Hmat \\
    \bm 0 & \bm 0 & \Hmat & \bm 0 & \Hmat' & -2\,\Hmat \\
    \bm 0 & \bm 0 & \bm 0 & \Hmat & 2\,\Hmat & \Hmat'
  \end{pmatrix}}_{\Amat \in\, \R^{6n\times 6L}}
  \bm\xi
  \;=\;
  \underbrace{\begin{pmatrix}
    -\bvec{C}''_x + 2\,\bvec{C}'_y + \tilde g_x \\
    -\bvec{C}''_y - 2\,\bvec{C}'_x + \tilde g_y \\
    \bm 0 \\ \bm 0 \\ \bm 0 \\ \bm 0
  \end{pmatrix}}_{\Bmat\,\in\,\R^{6n}},
\label{eq:linear-system}
\end{equation}

where the nominal-forcing terms that enter the first two blocks of
$\Bmat$ are

\begin{align}
  \tilde g_x(t) &= \Omegax(\bar x, \bar y)
                 - \Omegaxx(\bar x, \bar y)\,\bar x
                 - \Omegaxy(\bar x, \bar y)\,\bar y, \\
  \tilde g_y(t) &= \Omegay(\bar x, \bar y)
                 - \Omegaxy(\bar x, \bar y)\,\bar x
                 - \Omegayy(\bar x, \bar y)\,\bar y.
\end{align}

These are the ``constants of linearization'' --- what you have to
\emph{add} to the linearized dynamics to reproduce the full nonlinear
dynamics at the nominal state.

\begin{remark}
The rows are organised exactly as in Gino's rendezvous MATLAB code
\texttt{Rendezvous\_Gino.m}: two position-dynamics blocks, two
$\lamr$-dynamics blocks, two $\lamv$-dynamics blocks. The only
differences are (i) a planar problem rather than 3D, so the
block count drops from $9$ to $6$, and (ii) the HCW constant
matrices $M_{\mathrm{cw}}, N_{\mathrm{cw}}$ are replaced by the
\emph{time-varying} CR3BP Hessian entries evaluated on the nominal.
\end{remark}

\subsection{Number of unknowns and size of the linear system}

For $L$ ELM neurons per coordinate, $n$ training points, and the
planar problem:

\begin{center}
\begin{tabular}{lr}
  \toprule
  unknowns $\bm\xi$ & $6L$ \\
  equations         & $6n$ \\
  \bottomrule
\end{tabular}
\end{center}

We pick $n > L$ so the system is strictly overdetermined and
well-posed for least-squares solution. Typical choices for a first
Earth--Moon benchmark: $n = 40$, $L = 20$, giving a
$240\times 120$ linear system. The downstream Tikhonov+SVD step will
then reduce this to a few dozen ``effective'' unknowns before binary
encoding.


% ==================================================================
\section{Linear LSQ $\to$ QUBO pipeline}
% ==================================================================

This is the pipeline Gino introduced in \texttt{Rendezvous\_GIno.m}
and which the \texttt{qalunar.qubo.linear\_lsq} module ports to
Python. Given the linear system $\Amat \bm\xi = \Bmat$ from
Section~\ref{sec:linear-system}, construct a quadratic unconstrained
binary optimization (QUBO) whose ground state is the
Tikhonov-regularized least-squares solution under a per-variable
adaptive binary encoding.

\subsection{Step 1 --- Tikhonov regularization}

Minimize

\begin{equation}
  J(\bm\xi) \;=\; \|\Amat\bm\xi - \Bmat\|^2 \;+\; \lambda\,\|\bm\xi\|^2,
\label{eq:tikhonov}
\end{equation}

which is the standard regularized least-squares objective. This is
equivalent to the augmented system

\begin{equation}
  \Amat_{\mathrm{aug}} \;=\; \begin{pmatrix} \Amat \\ \sqrt\lambda \, \bvec{I}_{6L} \end{pmatrix},
  \qquad
  \Bmat_{\mathrm{aug}} \;=\; \begin{pmatrix} \Bmat \\ \bm 0 \end{pmatrix},
\end{equation}

and
\begin{equation}
  \min_{\bm\xi} \|\Amat_{\mathrm{aug}}\bm\xi - \Bmat_{\mathrm{aug}}\|^2.
\end{equation}

The regularization strength is chosen relative to the spectrum of
$\Amat$:

\begin{equation}
  \lambda \;=\; \varepsilon_{\mathrm{tik}} \cdot \sigma_{\max}(\Amat)^2,
\end{equation}

with $\varepsilon_{\mathrm{tik}} = 10^{-6}$ in the default config. The
scale $\sigma_{\max}^2$ makes the regularization self-calibrating to
the problem size; $\varepsilon_{\mathrm{tik}}$ controls how much the
regularizer bites into the solution.

\paragraph{Why bother with Tikhonov?}
Without it, rank-deficient or nearly-singular $\Amat$ (which is the
rule rather than the exception when $n \to L$ or when the ELM
spectrum decays fast) produces a solution with unboundedly large
entries $\bm\xi$. Tikhonov forces the minimum-norm-\emph{ish}
solution and gives a \emph{unique} $\bm\xi_{\mathrm{tik}}$ that we
can use as the encoding center.

\subsection{Step 2 --- Truncated SVD}

Compute the thin SVD

\begin{equation}
  \Amat_{\mathrm{aug}} \;=\; \bvec{U}\,\mathrm{diag}(\bm s)\,\bvec{V}^{\top},
\end{equation}

and keep only the singular values above a relative threshold:

\begin{equation}
  r \;=\; \#\{ i : s_i > \tau_{\mathrm{svd}}\, s_1 \},
  \qquad \tau_{\mathrm{svd}} = 10^{-3}\ \text{(default)}.
\end{equation}

The right-basis columns $\Vr \in \R^{6L \times r}$ define a
\emph{reduced} parameterization

\begin{equation}
  \bm\xi \;=\; \Vr\,\bm\alpha, \qquad \bm\alpha \in \R^{r},
\end{equation}

with the Tikhonov solution in reduced coordinates

\begin{equation}
  \bm\alpha_{\mathrm{tik}} \;=\; \mathrm{diag}(s_r)^{-1}
     \,\bvec{U}_r^{\top} \Bmat_{\mathrm{aug}},
\end{equation}

where $\bvec{U}_r$ and $s_r$ are the first $r$ columns and entries of
$\bvec{U}$ and $\bm s$. This typically cuts the number of variables
from hundreds to tens: the binary encoding only has to cover the
$r$-dimensional ``active'' subspace instead of the full $6L$ raw
weights.

\subsection{Step 3 --- Per-variable adaptive encoding}

Each reduced coordinate $\alpha_i$ is quantized over a personalized
interval $[\alpha_{\mathrm{tik},i} - R_i,\ \alpha_{\mathrm{tik},i} +
R_i]$ with $p$ bits. The half-range $R_i$ blends two contributions:

\begin{equation}
  R_i \;=\; \text{margin} \cdot
    \left(\,
      |\alpha_{\mathrm{tik},i}|
      \;+\;
      \frac{\|\bvec{r}_{\mathrm{aug}}\|}{s_i}
    \right),
\qquad
\bvec{r}_{\mathrm{aug}} = \Bmat_{\mathrm{aug}} - \Amat_{\mathrm{aug}}\bm\xi_{\mathrm{tik}},
\end{equation}

and is then clipped below by $\texttt{min\_spread}$ to avoid
degenerate zero-width intervals.

\paragraph{Why blend magnitude and sensitivity?}
The magnitude term $|\alpha_{\mathrm{tik},i}|$ reserves enough room
for the Tikhonov estimate itself (so the encoding isn't biased toward
zero on large components). The sensitivity term $\|r\| / s_i$
reserves more room for small-$s_i$ directions, which are the ones
that can move a lot without changing the LSQ residual much; it is
exactly the linear-sensitivity bound $|\delta \alpha_i| \leq \|\delta
b\| / s_i$ for a perturbed right-hand side.

The encoding is then

\begin{equation}
  \alpha_i(\bvec{q}) \;=\; \alpha_{\mathrm{offset},i}
     \;+\; \text{step}_i \sum_{k=0}^{p-1} 2^k\, q_{ip + k},
\qquad
\text{step}_i = \frac{2 R_i}{2^p - 1},
\quad
\alpha_{\mathrm{offset},i} = \alpha_{\mathrm{tik},i} - R_i.
\end{equation}

In matrix form, $\bm\alpha = \bm\alpha_{\mathrm{offset}} + \Ddec
\bvec{q}$, with $\Ddec \in \R^{r \times rp}$ block-diagonal with
blocks $\text{step}_i \cdot (1, 2, 4, \dots, 2^{p-1})$.

\subsection{Step 4 --- QUBO assembly}

Substituting the encoding into the Tikhonov objective:

\begin{align}
  J(\bvec{q})
   &\;=\; \|\Amat_{\mathrm{aug}} \Vr\,(\bm\alpha_{\mathrm{offset}} + \Ddec \bvec{q})
           - \Bmat_{\mathrm{aug}}\|^2 \\
   &\;=\; \|\Phi\,\bvec{q} - \tilde{\Bmat}\|^2,
\end{align}

where

\begin{align}
  \Amat_{\mathrm{red}} &= \Amat_{\mathrm{aug}}\, \Vr, \\
  \Phi                 &= \Amat_{\mathrm{red}}\,\Ddec, \\
  \tilde{\Bmat}        &= \Bmat_{\mathrm{aug}} - \Amat_{\mathrm{red}}\,\bm\alpha_{\mathrm{offset}}.
\end{align}

Expanding the square gives the QUBO

\begin{equation}
\boxed{
  E(\bvec{q}) \;=\; \bvec{q}^{\top} \Qmat \bvec{q}
              \;+\; \bvec\ell^{\top}\bvec{q}
              \;+\; c,
}
\end{equation}

with

\begin{equation}
  \Qmat = \Phi^{\top}\Phi,\qquad
  \bvec\ell = -2\,\Phi^{\top}\tilde{\Bmat},\qquad
  c = \|\tilde{\Bmat}\|^2.
\end{equation}

\paragraph{The core identity.}
By construction,
\begin{equation}
  E(\bvec{q}) \;=\; \|\Amat_{\mathrm{aug}} \bm\xi(\bvec{q}) - \Bmat_{\mathrm{aug}}\|^2
  \quad\text{for every }\bvec{q}\in\{0,1\}^{rp}.
\end{equation}
That is, the QUBO ground state \emph{is} the best binary-encodable
approximation to the Tikhonov-regularized LSQ solution --- no
approximation error beyond quantization. The unit tests in
\texttt{tests/test\_qubo\_linear\_lsq.py} verify this identity on
random bit vectors.

\subsection{Decoding and warm starts}

To go from a bitstring $\bvec{q}$ back to the original unknowns:

\begin{equation}
  \bvec{q} \;\;\longrightarrow\;\;
  \bm\alpha = \bm\alpha_{\mathrm{offset}} + \Ddec\,\bvec{q}
  \;\;\longrightarrow\;\;
  \bm\xi = \Vr\,\bm\alpha.
\end{equation}

The \textbf{warm start} is the quantization of $\bm\alpha_{\mathrm{tik}}$
to $p$ bits per coordinate. Because the encoding is centered at
$\bm\alpha_{\mathrm{tik}}$, the warm start is within half a quantization
step per coordinate of the Tikhonov solution, and the corresponding
bitstring is a strong initial seed for D-Wave hybrid solvers or
simulated annealing.

\subsection{Default hyperparameters}

\begin{center}
\begin{tabular}{lll}
  \toprule
  Parameter & Symbol & Default \\
  \midrule
  Bits per reduced variable        & $p$                     & $8$     \\
  Tikhonov (relative to $\sigma_{\max}^2$) & $\varepsilon_{\mathrm{tik}}$ & $10^{-6}$ \\
  SVD truncation threshold         & $\tau_{\mathrm{svd}}$   & $10^{-3}$ \\
  Encoding margin                  & --                      & $3.0$   \\
  Minimum half-range floor         & --                      & $10^{-6}$ \\
  \bottomrule
\end{tabular}
\end{center}

These match Gino's MATLAB defaults and are the starting point for
the sensitivity sweep once the Earth--Moon benchmark runs.

% ==================================================================
\section{Sequential linearization outer loop}\label{sec:sequential-linearization}
% ==================================================================

Because the CR3BP is nonlinear, one call to the QUBO inner solver
does not give the optimal trajectory, it gives the optimal
\emph{linear perturbation} around whatever nominal was used to build
$\Amat, \Bmat$. The sequential outer loop iterates:

\begin{enumerate}[leftmargin=1.5em]
  \item Start with an initial nominal $\bar{\xvec}^{(0)}(t)$
    (e.g.\ the Hermite cubic interpolant of the boundary conditions,
    or a classical-NLP reference solution).
  \item At iteration $k$:
    \begin{enumerate}[leftmargin=1.5em, itemsep=1pt]
      \item Evaluate $\Omegaxx, \Omegaxy, \Omegayy$ and the forcing
        terms $\tilde g_x, \tilde g_y$ on the nominal.
      \item Assemble the linear system $\Amat^{(k)} \bm\xi = \Bmat^{(k)}$.
      \item Build and solve the QUBO; decode to $\bm\xi^{(k)}$ and
        reconstruct the new trajectory $\xvec^{(k+1)}(t)$.
      \item Update the nominal: $\bar{\xvec}^{(k+1)} = \xvec^{(k+1)}$.
    \end{enumerate}
  \item Repeat until $\|\bar{\xvec}^{(k+1)} - \bar{\xvec}^{(k)}\|$
    drops below a tolerance, or the cost $J$ stops decreasing.
\end{enumerate}

This is essentially a Gauss--Newton / sequential linearization
approach, very similar to what sequential convex programming (SCP)
does in the classical trajectory-optimization literature. The
difference is that the inner step is not a convex QP solve but a
QUBO ground-state search.

\paragraph{Pass-1 vs.\ Pass-2 loop.}
The implementation in
\texttt{IndirectTfcElmTranscription.solve\_sequential} supports two
modes, selected by the flag \texttt{OuterLoopConfig.line\_search}:

\begin{itemize}[leftmargin=1.5em, itemsep=1pt]
  \item \textbf{Pass 1} (\texttt{line\_search=False}): fixed damping,
    $\bar{\xvec}^{(k+1)} = \bar{\xvec}^{(k)} + \alpha\,(\xvec^{(k+1)} - \bar{\xvec}^{(k)})$
    with a constant $\alpha \in (0, 1]$. Reproduces the behaviour of
    the first-pass rendezvous loop. Tolerant of mild nonlinearity;
    oscillates or diverges on stronger CR3BP transfers.
  \item \textbf{Pass 2} (\texttt{line\_search=True}, the default):
    adaptive Armijo backtracking with merit function
    $m(\bar{\xvec}) = \|\xvec(\bm\xi(\bar{\xvec})) - \bar{\xvec}\|_\infty$,
    i.e.\ the \emph{fixed-point residual} of the Picard map
    $\bar{\xvec} \mapsto \xvec(\bm\xi(\bar{\xvec}))$, which is the
    right Newton merit for the rootfinding problem
    $F(\bar{\xvec}) := \xvec(\bm\xi(\bar{\xvec})) - \bar{\xvec} = \bm 0$.
\end{itemize}

\subsection{Pass-2 Armijo line search}
\label{sec:armijo}

The full step at iteration $k$ is
\[
  \Delta \bar{\xvec}^{(k)} \;=\; \xvec^{(k+1)} - \bar{\xvec}^{(k)}.
\]
Starting from $\alpha = \texttt{damping}$ (default $1.0$ --- a full
Newton-like step), we trial
\[
  \bar{\xvec}^{(k+1)}_{\mathrm{trial}} \;=\; \bar{\xvec}^{(k)}
     \;+\; \alpha \, \Delta \bar{\xvec}^{(k)}
\]
and accept it when
\begin{equation}
  m\big(\bar{\xvec}^{(k+1)}_{\mathrm{trial}}\big)
  \;\leq\;
  (1 - c_1\,\alpha)\, m\big(\bar{\xvec}^{(k)}\big),
\label{eq:armijo}
\end{equation}
with $c_1 = \texttt{armijo\_c} = 10^{-4}$. Otherwise $\alpha$ is
multiplied by \texttt{backtrack\_factor} (default $1/2$) and the
trial is repeated, down to a floor $\alpha \geq \texttt{min\_damping}$
(default $2^{-10}$).

\paragraph{Why the fixed-point residual?}
The nonlinear TPBVP residual is \emph{not} a valid Armijo merit
here: the linearization changes at every trial nominal, so
$\|R_{\mathrm{NL}}(\bm\xi^{\mathrm{trial}})\|$ need not be monotone
with $\alpha$. The fixed-point residual
$\|F(\bar{\xvec})\|_\infty$, on the other hand, is exactly the
quantity the outer loop is trying to zero; it is the analogue of
$\|\nabla\!f\|$ in gradient descent, and its Armijo-like decrease
gives the outer loop a globalization guarantee --- once Pass~2 has
accepted a step, the fixed-point residual cannot grow above its
initial value.

\paragraph{Concrete effect.}
On the short, gentle fixture used in
\texttt{tests/test\_indirect\_tfc\_elm.py}, Pass~2 accepts
$\alpha = 1$ at every iteration and converges in about $13$ outer
iterations, versus $\sim\! 32$ for fixed damping $\alpha = 0.5$. On
the long-transfer fixture, fixed-damping $\alpha = 1$ oscillates
wildly --- peak step norm $\approx 2$ against an initial value
$\approx 0.7$ --- while the Armijo-bounded version keeps the step
norm under its initial value for all $40$ iterations, as the test
\texttt{TestLineSearch::test\_line\_search\_bounds\_step\_norm\_on\_%
aggressive\_problem} verifies. That said, bounding the step norm is
not the same as converging: on the long transfer, even the
line-searched loop stalls outside the basin of attraction if seeded
only by the BC-only Hermite interpolant, which motivates the
classical reference of the next section.

\newpage
% ==================================================================
\section{Classical reference via direct collocation}
\label{sec:reference}
% ==================================================================

The outer loop of Section~6 is a contraction \emph{only inside a
local basin of attraction}, and the basin can be genuinely small on
long or curved transfers in the Earth--Moon CR3BP. The indirect
transcription needs a good initial nominal, and the BC-only Hermite
interpolant \texttt{initial\_nominal()} is not good enough once the
trajectory makes any real excursion away from the straight line
joining the endpoints.

We therefore ship a \textbf{classical direct-collocation reference
solver} in \texttt{qalunar.reference.direct\_collocation}. It plays
two distinct roles in the project:

\begin{enumerate}[leftmargin=1.5em]
  \item \textbf{Warm start.} The solution produced by the reference
    solver is resampled onto the TFC+ELM training grid and passed as
    \texttt{initial\_nominal} to
    \texttt{solve\_sequential}.
  \item \textbf{Classical baseline.} The same solver provides a CPU
    reference objective $J^{\star}_{\mathrm{NLP}}$ that the QUBO
    samples will be measured against in the benchmark sweep
    (project Task~9).
\end{enumerate}

\subsection{Hermite--Simpson direct collocation}

Split $[0, T]$ into $N$ uniform intervals of width $h = T/N$ and
introduce node values at $N+1$ points,
\[
  \xvec_k \;=\; (x_k, y_k, v_{x,k}, v_{y,k}),
  \qquad
  \uvec_k \;=\; (u_{x,k}, u_{y,k}),
  \qquad k = 0, \dots, N.
\]
The decision vector is the concatenation
\[
  \bvec{z} \;=\; (\xvec_0, \dots, \xvec_N,\;\uvec_0, \dots, \uvec_N)
     \;\in\; \R^{6(N+1)}.
\]

\paragraph{Hermite--Simpson defect.}
On each sub-interval $[t_k, t_{k+1}]$, the cubic Hermite interpolant
through $(\xvec_k, \fvec_k)$ and $(\xvec_{k+1}, \fvec_{k+1})$ must
agree with the dynamics at the midpoint. With
$\fvec_k = \fvec(\xvec_k, \uvec_k)$,
\begin{align}
  \xvec_{k+1/2} &= \tfrac12(\xvec_k + \xvec_{k+1})
     \;+\; \tfrac{h}{8}\,(\fvec_k - \fvec_{k+1}), \\
  \uvec_{k+1/2} &= \tfrac12(\uvec_k + \uvec_{k+1}), \\
  \fvec_{k+1/2} &= \fvec(\xvec_{k+1/2}, \uvec_{k+1/2}).
\end{align}
The Hermite--Simpson defect constraint is
\begin{equation}
  \boxed{\;
  \bvec{d}_k \;=\; \xvec_{k+1} - \xvec_k
     \;-\; \frac{h}{6}\big(\fvec_k + 4\,\fvec_{k+1/2} + \fvec_{k+1}\big)
     \;=\; \bm 0,
  \;}
  \qquad k = 0, \dots, N - 1.
\label{eq:hs-defect}
\end{equation}
This is a third-order collocation scheme, and it is a standard
choice from Kelly~\cite{Kelly2017} and Betts~\cite{Betts2010}.

\paragraph{Objective.}
The energy cost is discretized with the composite trapezoidal rule,
\begin{equation}
  \tilde J(\bvec z) \;=\; \frac{h}{2}\sum_{k=0}^{N} w_k\,\|\uvec_k\|^2,
  \qquad w_0 = w_N = \tfrac12,
  \qquad w_k = 1\ \text{otherwise}.
\end{equation}
$\tilde J$ is a pure quadratic in the control block (Hessian constant
and sparse), and its gradient
\[
  \partial \tilde J/\partial \uvec_k \;=\; h\,w_k\,\uvec_k
\]
is handed analytically to SLSQP; only the defect Jacobian is
finite-differenced.

\paragraph{Boundary conditions and bounds.}
The eight equalities
$\xvec_0 = (\bvec r_0, \bvec v_0)$ and
$\xvec_N = (\bvec r_f, \bvec v_f)$ are appended to the constraint
list. An optional symmetric box
$|u_{x,k}|, |u_{y,k}| \leq u_{\max}$ can be passed via
\texttt{DirectCollocationConfig.control\_bound}; states are left
unbounded.

The full NLP is
\begin{equation}
  \min_{\bvec z}\; \tilde J(\bvec z)
  \quad \text{s.t.}\quad
  \bvec d_k(\bvec z) = \bm 0\ \forall k,\ \
  \xvec_0 = \bvec s_0,\ \
  \xvec_N = \bvec s_f,\ \
  \uvec_k \in [-u_{\max}, u_{\max}]^2,
\end{equation}
and is handed to \texttt{scipy.optimize.minimize} with
\texttt{method="SLSQP"}. With $N = 40$ the problem has a few hundred
variables and constraints and solves in well under a second.

\subsection{Resampling onto the TFC+ELM grid}

The direct-NLP result is returned as
\texttt{DirectCollocationResult}, which exposes a helper
\begin{center}
\texttt{sample\_position(t\_query) -> (x\_bar, y\_bar)}
\end{center}
that linearly interpolates the node positions onto an arbitrary time
grid. Warm-starting the indirect solver is then a three-liner:
\begin{lstlisting}[basicstyle=\ttfamily\small]
ref = solve_energy_optimal_cr3bp(dynamics, r0, v0, rf, vf, T)
x_bar, y_bar = ref.sample_position(transcription.t)
sol = transcription.solve_sequential(
    OuterLoopConfig(), initial_nominal=(x_bar, y_bar),
)
\end{lstlisting}
Linear interpolation is enough because the outer loop uses the
nominal only to evaluate $\Omegaxx, \Omegaxy, \Omegayy$ and
$\tilde g_x, \tilde g_y$ pointwise; those are smooth functions of
$(x, y)$ and insensitive to sub-node interpolation error.

\subsection{Example trajectories: a two-phase Earth--Moon mission}
\label{sec:example-trajectories}

A physically realistic LEO-to-lunar low-thrust transfer is a
multi-revolution spiral with hundreds of orbits around Earth over
months to years of flight, well beyond what a single polynomial (or
ELM) expansion can represent. Rather than ask the CR3BP
transcription to do work it cannot do, we split the mission in two:

\begin{itemize}[leftmargin=1.5em, itemsep=1pt]
  \item \textbf{Phase 1} --- analytical Edelbaum spiral from a low
    circular Earth orbit out to a high circular Earth orbit, ignoring
    lunar gravity. In this regime the dynamics are pure two-body
    Keplerian around Earth and the problem is closed-form: tangential
    thrust is already Edelbaum-optimal for a coplanar circular-to-
    circular transfer, so there is nothing to optimize. We read off
    $\Delta v$, time-of-flight and revolution count from the
    classical formulas and move on.
  \item \textbf{Phase 2} --- CR3BP energy-optimal transfer from the
    Phase-1 handoff state to a rest rendezvous west of the Moon,
    solved by the Hermite--Simpson direct NLP described above (and,
    later, by the TFC+ELM indirect transcription of
    Section~\ref{sec:transcription}). This is the three-body leg
    that the whole QUBO stack targets.
\end{itemize}

Phase 1 is implemented in \texttt{qalunar.reference.edelbaum} as a
ten-line closed-form function
\texttt{edelbaum\_spiral}$(r_1, r_2, a_t)$ that returns an
\texttt{EdelbaumSpiral} dataclass. In CR3BP nondimensional units,
with $GM_{\oplus} = 1 - \mu$, the inertial circular speed is
$v_i = \sqrt{(1-\mu)/r_i}$, so the total $\Delta v$ is just
$|v_1 - v_2|$, the time of flight is $|v_1 - v_2| / a_t$ (constant
tangential thrust integrates trivially), and the revolution count
comes from integrating the orbital angular rate $v/r = v^3/GM$
against $dt$:
\begin{equation}
  N_{\mathrm{rev}}
  \;=\;
  \frac{|v_1^4 - v_2^4|}{\,8\pi\,(1-\mu)\,a_t\,}.
\end{equation}
The \texttt{handoff\_state\_synodic()} method then places the
spacecraft on the $+x$ direction at the moment its Earth-relative
position crosses the positive $x$-axis, so its synodic position is
$(-\mu + r_2, 0)$ with inertial velocity $(0, v_2)$ prograde, mapped
to synodic velocity by subtracting
$\bm\omega \times \bvec r = (-y, x)$.

\paragraph{The reference mission.}
The demo hero mission is a spiral from LEO ($r_1 = 0.018$ nondim,
$\sim\!6\,900$ km) to a high circular Earth orbit at $r_2 = 0.5$ nondim
($\sim\!192\,000$ km, half the Earth--Moon distance), at a plausible
solar-electric thrust acceleration of $0.1\,\mathrm{mm/s^2}$. Phase 1
yields (all numbers from \texttt{EdelbaumSpiral.si\_report()}):
\begin{itemize}[leftmargin=1.5em, itemsep=1pt]
  \item $\Delta v_{\mathrm{P1}} \approx 6\,142$ m/s,
  \item $T_{\mathrm{P1}} \approx 709$ days $\approx 1.94$ yr,
  \item $N_{\mathrm{rev}} \approx 3\,292$ revolutions around Earth.
\end{itemize}
The handoff state is
$\bvec r_0 = (-\mu + 0.5, 0) \approx (0.488, 0)$ with synodic velocity
$\bvec v_0 = (0,\, v_2 - (-\mu + r_2)) \approx (0, 0.918)$ --- a
prograde tangential crossing of the $+x$ axis on the Phase-1 exit
orbit. Phase 2 then runs the direct NLP from this state to a rest
rendezvous at $\bvec r_f = (1-\mu-0.1, 0)$,
$\bvec v_f = \bvec 0$, with flight time $T = 3.0$ nondim
($\approx 13$ days in physical units). SLSQP drives the
Hermite--Simpson defects to $\approx 3 \cdot 10^{-10}$ in $53$
iterations with objective
$J = \tfrac12 \int_0^T \|\uvec\|^2\, dt \approx 0.64$ and peak
acceleration $\|\uvec\|_{\max} \approx 1.69$ nondim. The resulting
arc is a single prograde sweep from the Phase-1 exit orbit to the
Moon realm, shown in Figure~\ref{fig:reference-gentle}.

\begin{figure}[htbp]
  \centering
  \includegraphics[width=0.98\textwidth]{../scripts/reference_gentle.png}
  \caption{Two-phase Earth--Moon low-thrust mission. Phase 1 (dashed
    gray) is an analytical Edelbaum spiral from LEO out to a
    $r = 0.5$ nondim circular Earth orbit, summarized in the caption
    box; the spiral itself is hundreds of revolutions and is not
    drawn explicitly. Phase 2 (orange) is the direct-NLP
    energy-optimal CR3BP transfer from the Phase-1 handoff state
    $(0.488, 0)$ with synodic velocity $(0, 0.918)$ to a rest
    rendezvous at $(1-\mu-0.1, 0)$, flight time $T = 3.0$ nondim.
    Left: position in the synodic frame, with Earth and Moon
    markers. Right: Phase-2 thrust-acceleration magnitude, showing a
    smooth push-brake profile with peak $\approx 1.69$ nondim and
    objective $J \approx 0.64$. The Phase-1 summary box reports
    $\Delta v = 6\,142$ m/s, $T = 709$ days, $N_{\mathrm{rev}} =
    3\,292$ at $0.1$ mm/s$^2$.}
  \label{fig:reference-gentle}
\end{figure}

A second, deliberately synthetic benchmark (the same BCs used by
the \texttt{transcription} test fixture: $\bvec r_0 = (-0.3, 0)$,
$\bvec v_0 = (0, 0.6)$, $\bvec r_f = (0.4, 0.2)$,
$\bvec v_f = (-0.1, 0)$, $T = 2.5$) is what exhibits the warm-start
advantage. These BCs were chosen to keep the trajectory \emph{away}
from both primaries so that the Hessian entries $\Omegaxx, \Omegaxy,
\Omegayy$ (which scale as $\sim\! 1/r^3$ near each primary) stay
bounded and unit-test conditioning is predictable --- they are a
numerical stress test for the indirect iteration, not a mission
geometry. On the two-phase Earth--Moon mission of
Figure~\ref{fig:reference-gentle}, direct NLP and the indirect
TPBVP sit in nearby basins and cold-start already contracts almost
as well as warm-start, so the warm-start \emph{advantage} is best
displayed on the synthetic benchmark, where cold-start is
catastrophically stuck.
The synthetic-benchmark comparison overlays, in the synodic frame:

\begin{itemize}[leftmargin=1.5em, itemsep=1pt]
  \item the direct-NLP reference (orange),
  \item the cold-started indirect iteration --- BC-only Hermite seed,
    no warm start --- after $40$ outer iterations (red dashed),
  \item the indirect iteration warm-started from
    \texttt{ref.sample\_position(transcription.t)} (green).
\end{itemize}

The cold-started indirect trajectory veers off into a completely
unphysical arc: the sequential linearization is stuck outside the
basin of attraction and never recovers, exactly the failure mode the
Pass-2 Armijo line search \emph{bounds} but does not cure. The
warm-started iteration, in contrast, tracks the classical reference
within a handful of outer iterations.

\IfFileExists{../scripts/warm_start_aggressive.png}{%
\begin{figure}[htbp]
  \centering
  \includegraphics[width=0.72\textwidth]{../scripts/warm_start_aggressive.png}
  \caption{Synthetic long-transfer stress test in the synodic frame,
    three overlays: the direct-NLP classical reference (orange), the
    cold-started indirect iteration (red dashed), and the same
    indirect iteration warm-started from the reference (green). Only
    the warm-started solve tracks the reference; the cold-started
    iteration stalls outside the basin of attraction. These BCs are
    chosen to keep the trajectory away from both primaries so the
    $\Omega$-Hessian stays bounded --- they are a numerical benchmark
    for the indirect iteration, not a mission geometry.}
  \label{fig:warm-start}
\end{figure}
}{}

Quantitatively, the nonlinear TPBVP residual history for the two
seeds tracks the same pattern. The BC-only seed plateaus
at $\|R_{\mathrm{NL}}\|_\infty \approx 13$ after two iterations and
never moves. The direct-NLP warm start drops from $\approx 30$ to
$\approx 0.74$ in three iterations --- roughly a $17\times$
reduction --- and plateaus there. The remaining $\sim\! 0.74$
residual is the last obstacle: it is a fixed point of the
\emph{linearization} whose $\Omega$-Hessian freezing drift is small
enough to be in the basin of the true TPBVP but too large to be
zeroed by the pure Picard structure of the outer loop, which only
relinearizes at committed nominals and therefore carries stale
$A, B$ matrices inside each Armijo trial step. Cracking this
remainder motivates a Levenberg--Marquardt / Gauss--Newton extension
of \texttt{solve\_sequential} that relinearizes at every trial step
during the backtrack; a direct-LGL pseudospectral transcription is
held as a backup baseline for comparison. Project Task~8 will
report whether the $\sim\! 0.74$ remainder matters for QUBO
quality in practice.

\IfFileExists{../scripts/convergence_history.png}{%
\begin{figure}[htbp]
  \centering
  \includegraphics[width=0.85\textwidth]{../scripts/convergence_history.png}
  \caption{Nonlinear TPBVP residual vs.\ outer iteration on the
    synthetic long-transfer benchmark, cold-start vs.\ warm-start
    (log scale on the $y$-axis). The warm start drops from
    $\approx 30$ to $\approx 0.74$ in three iterations; the cold
    start plateaus near $13$ and does not move.}
  \label{fig:convergence}
\end{figure}
}{}

\newpage
% ==================================================================
\section{Code organization}
% ==================================================================

The codebase lives under \texttt{qalunar/} and is organized by
separation of concerns:

\begin{lstlisting}[basicstyle=\ttfamily\small]
qalunar/
  __init__.py
  dynamics/
    __init__.py
    cr3bp.py                  # PlanarCR3BP: rhs, Jacobian, Jacobi integral
  qubo/
    __init__.py
    linear_lsq.py             # build_linear_lsq_qubo + LinearLsqToQuboResult
  transcription/
    __init__.py
    indirect_tfc_elm.py       # IndirectTfcElmTranscription + OuterLoopConfig
  reference/
    __init__.py
    direct_collocation.py     # Hermite-Simpson NLP (warm start + baseline)
    edelbaum.py               # analytical Phase-1 Earth-escape spiral
tests/
  test_cr3bp.py                        # 18 tests: L4/L5, df/dx vs FD, ...
  test_qubo_linear_lsq.py              # 24 tests: energy identity, warm start
  test_indirect_tfc_elm.py             # 51 tests: TFC, ELM, outer loop, line search
  test_reference_direct_collocation.py # 20 tests: HS defects, warm-start
  test_reference_edelbaum.py           # 26 tests: closed-form spiral, handoff
scripts/
  plot_cr3bp_overview.py              # primaries, L4/L5, ZVCs
  plot_reference_trajectories.py      # Section 7 figures
docs/
  qalunar_manual.tex          # this document
\end{lstlisting}

\subsection{Running the tests}

\begin{lstlisting}[language=bash]
# from the project root
python -m pytest -q
\end{lstlisting}

The full suite currently contains 139 tests and runs in roughly 2
minutes (the bulk of the runtime is the SLSQP solves inside the
direct-collocation tests). It is the first place to look when
something is misbehaving: every mathematical identity in this manual
corresponds to at least one executable assertion.

\subsection{Key modules at a glance}

\paragraph{\texttt{qalunar.dynamics.cr3bp}.}
Planar CR3BP with a 2D thrust control. \texttt{PlanarCR3BP.rhs} and
\texttt{jacobian\_state} implement equations
\eqref{eq:cr3bp-1st-order} and the state Jacobian used for
linearization. A \texttt{lagrange\_triangular\_points} helper
returns $L_4$, $L_5$ symbolically. The Earth--Moon mass parameter is
exposed as the module constant \texttt{EARTH\_MOON\_MU}.

\paragraph{\texttt{qalunar.qubo.linear\_lsq}.}
One-function entry point:
\texttt{build\_linear\_lsq\_qubo(A, B, config)}. Returns a
\texttt{LinearLsqToQuboResult} bundling the QUBO
$(\Qmat, \bvec\ell, c)$, the decoding path $(\Vr,
\bm\alpha_{\mathrm{offset}}, \text{step}_i)$, and diagnostics
(singular spectrum, effective rank, Tikhonov solution,
$\lambda_{\mathrm{tik}}$). Methods:
\texttt{decode\_alpha}, \texttt{decode\_xi}, \texttt{energy},
\texttt{warm\_start\_bits}.

\paragraph{\texttt{qalunar.transcription.indirect\_tfc\_elm}.}
The indirect TFC+ELM pipeline of Sections 4--6.
\texttt{IndirectTfcElmTranscription} assembles the linear system of
\eqref{eq:linear-system} for a given nominal via
\texttt{build\_linear\_system}, and
\texttt{solve\_sequential(cfg, initial\_nominal=...)} runs the outer
loop: classical LSQ solve, decode, Pass-2 Armijo backtracking
(Section~\ref{sec:armijo}), update nominal. Returns a
\texttt{SequentialSolution} with \texttt{xi}, \texttt{trajectory} and
a per-iteration \texttt{history} dict containing
\texttt{trajectory\_change\_inf}, \texttt{linearized\_residual},
\texttt{nonlinear\_residual} and \texttt{step\_size}.

\paragraph{\texttt{qalunar.reference.direct\_collocation}.}
The classical Hermite--Simpson direct collocation solver of
Section~\ref{sec:reference}. Entry points:
\texttt{DirectCollocationConfig}, \texttt{DirectCollocationResult}
and the top-level function
\texttt{solve\_energy\_optimal\_cr3bp(dynamics, r0, v0, rf, vf, T,
config)}. The returned result carries \texttt{t, x, y, vx, vy, ux, uy},
the SLSQP convergence status, and a \texttt{sample\_position(t\_query)}
helper used to warm-start
\texttt{IndirectTfcElmTranscription.solve\_sequential}.

\paragraph{\texttt{qalunar.reference.edelbaum}.}
Closed-form analytical solver for the Phase-1 Earth-escape spiral
(Section~\ref{sec:example-trajectories}).
\texttt{edelbaum\_spiral(r1, r2, a\_thrust, mu)} returns an
\texttt{EdelbaumSpiral} dataclass holding $r_i, v_i, \Delta v, T,
N_{\mathrm{rev}}, a_t$ in CR3BP nondimensional units. The
\texttt{handoff\_state\_synodic()} method yields the synodic
$(\bvec r, \bvec v)$ initial state for the Phase-2 direct NLP, and
\texttt{si\_report()} returns a dict of pre-formatted SI-unit strings
($\Delta v$ in m/s, TOF in days, etc.) for manual captions and
progress logs.

The indirect TFC+ELM pipeline produces a linear system
$\Amat \bm\xi = \Bmat$ with decoding metadata compatible with
\texttt{qubo.linear\_lsq}; a direct-LGL pipeline will land alongside
it in \texttt{qalunar.transcription} as part of project Task~8.


% ==================================================================
\section{Convergence investigation and pivot}\label{sec:convergence-investigation}
% ==================================================================

The sequential linearization loop of the hybrid QUBO--classical pipeline
was expected to converge toward the known direct-NLP solution.
Extensive numerical experiments revealed two orthogonal failure modes
that together make the continuous-lstsq QUBO approach unsuitable for
production-quality trajectory optimization.

\subsection{Armijo line search failure}

The Armijo backtracking strategy (Pass~2) uses the fixed-point residual
$m(\bar{\xvec}) = \|\xvec(\bm\xi(\bar{\xvec})) - \bar{\xvec}\|_\infty$
as merit function.
Diagnostic runs showed that this merit function is too noisy for the
Picard map in the CR3BP: the Armijo condition is never satisfied at
large step sizes, forcing $\alpha$ down to the floor
$\alpha_{\min} = 0.001$ at every iteration.
The effective update becomes negligibly small, and the nonlinear
residual stagnates rather than converging.

In contrast, \emph{undamped Picard iteration} (fixed $\alpha=1$, no
line search) converges reliably on the same problems.  Using the
classical NLP solution as an initial nominal and solving the inner
lstsq problem exactly (not via QUBO), undamped Picard reduces the
nonlinear residual from $5.7$ to $1.9 \times 10^{-5}$ in fewer than
10~iterations and $0.2$~seconds of wall time.

\begin{remark}
This confirms that the sequential linearization itself is sound; the
inner linear solve and the outer Picard update converge.  The Armijo
mechanism was the bottleneck, not the linearization quality.
\end{remark}

\subsection{QUBO quantization error}\label{sec:qubo-quantization}

Even after removing the Armijo failure (using $\alpha=1$), replacing
the exact lstsq inner solve with the QUBO pipeline introduces a
discretization error that prevents convergence.

Each continuous reduced variable $\alpha_j$ is encoded as a
$p$-bit binary expansion (Section~5):
\[
  \alpha_j \;\approx\; \alpha_j^{\min} + \Delta_j \sum_{k=0}^{p-1}
  2^k\, q_{jp+k},
  \qquad
  \Delta_j = \frac{\alpha_j^{\max} - \alpha_j^{\min}}{2^p - 1}.
\]
With $p=8$ bits, each variable has a resolution of roughly
$\Delta_j / 256$.  The QUBO ground state finds the best
\emph{representable} $\bm\alpha$, which generically differs from the
exact lstsq minimizer by $O(\Delta_j)$.

At Picard iteration $k$, the update reconstructed from the QUBO
solution carries this quantization error forward into the next
nominal.  Because the linearization point has shifted by a
quantization-scale perturbation rather than the true Newton step,
subsequent iterations cannot reduce the residual below a floor set by
$\Delta_j$.  Numerically, the hybrid loop with 8-bit encoding
oscillates between nonlinear residuals of $\sim 7$--$12$ indefinitely,
never approaching the $10^{-5}$ achievable by exact lstsq.

Increasing $p$ would shrink $\Delta_j$ but expands the QUBO dimension
(number of qubits) proportionally.  At $p=16$, the QUBO for a problem
with $r = 30$ reduced variables already requires $480$~qubits and a
$480 \times 480$ coupling matrix, and even 16-bit resolution
($\sim 10^{-5}$ per variable) may not suffice for demanding transfers.

\subsection{Why polishing does not save the approach}

One could ``polish'' the QUBO output by running a classical lstsq step
on the decoded continuous solution.  But this is scientifically
misleading: the lstsq polish does the actual work of finding the
continuous minimum, while the QUBO merely provides a warm start that
the lstsq solver does not need.  Classical lstsq solves the inner
problem in $\sim 5$~ms; the QUBO \emph{approximates} the same answer
in orders of magnitude more time.

\subsection{Conclusion: the continuous-lstsq QUBO is the wrong problem}

The Carbone et~al.\ approach~\cite{Carbone2025} works for the
particular problem it was designed for (HCW rendezvous with a single
linear solve), but does not generalize well to nonlinear problems
requiring iterative refinement.  The fundamental mismatch is that QUBO
variables are binary, yet the underlying optimization variables are
continuous.  Artificially encoding continuous variables as bitstrings
introduces quantization error that dominates the solution quality.

This motivates a search for problems where the decision variables are
\emph{inherently binary}, so that the QUBO formulation is exact rather
than an approximation.


% ==================================================================
\section{On/off thrust scheduling via QUBO}\label{sec:thrust-scheduling}
% ==================================================================

Low-thrust engines can operate in two modes: on (full thrust) or off
(coasting).  Given a transfer window divided into $N$ time steps, the
decision ``fire at step $i$?''\ is a naturally binary variable
$q_i \in \{0, 1\}$.  The search space of $2^N$ possible schedules is
combinatorially hard for classical exhaustive search, making this an
ideal candidate for quantum annealing.

\subsection{Problem setup}

Consider a spacecraft in the planar CR3BP with initial state
$\xvec_0 = (x_0, y_0, \dot x_0, \dot y_0)\TT$.
The time interval $[t_0, t_f]$ is divided into $N$ decision intervals
of length $\Delta t = (t_f - t_0)/N$.  At each interval $i$, the
engine is either on ($q_i = 1$, applying thrust acceleration $\uvec_i$)
or off ($q_i = 0$, coasting).

The thrust direction at each step can be:
\begin{itemize}[leftmargin=1.5em, itemsep=1pt]
  \item \textbf{tangential} --- along the velocity vector,
  \item \textbf{anti-tangential} --- opposite to the velocity vector, or
  \item \textbf{fixed} --- a constant direction $\hat{\uvec}$ specified
    by the user.
\end{itemize}
In all cases, the magnitude $\|\uvec_i\| = a_T$ is constant (on/off
model: full thrust or zero).

\subsection{Linearization via the state transition matrix}

To formulate a QUBO, the effect of each burn on the final state must
be expressed as a linear function of the binary variables $q_i$.

\paragraph{Coast arc and STM.}
First, propagate the uncontrolled (coast) trajectory from $\xvec_0$
over $[t_0, t_f]$ using the CR3BP equations of motion.  Simultaneously
integrate the state transition matrix (STM) via the variational
equations:
\[
  \dot{\Phi}(t, t_0) = \bvec A(t)\,\Phi(t, t_0),
  \qquad \Phi(t_0, t_0) = \bvec I_{4\times 4},
\]
where $\bvec A(t)$ is the Jacobian of the CR3BP right-hand side
evaluated along the coast arc.  The implementation integrates the
augmented 20-dimensional system
$\bvec z = (\xvec, \mathrm{vec}(\Phi))$ using fixed-step RK4.

The STM $\Phi(t_i, t_f)$ from decision point $t_i$ to the final time
is obtained by composition:
\[
  \Phi(t_i, t_f) = \Phi(t_f, t_0)\,\Phi(t_i, t_0)^{-1}.
\]

\paragraph{Impulse response vectors.}
The control influence matrix for the planar CR3BP is
\[
  \bvec B_{\mathrm{ctrl}} =
  \begin{pmatrix} 0 & 0 \\ 0 & 0 \\ 1 & 0 \\ 0 & 1 \end{pmatrix},
\]
so an impulsive thrust acceleration $\uvec_i$ at time $t_i$ produces
a first-order change at $t_f$ of:
\[
  \bvec b_i \;=\; \Phi(t_i, t_f)\, \bvec B_{\mathrm{ctrl}}\, \uvec_i\,
  \Delta t.
\]
The vector $\bvec b_i \in \R^4$ encodes the full effect of burning at
step $i$ on the final position and velocity.

\paragraph{Total effect of a schedule.}
The superposition principle (valid under the linearization) gives the
final-state deviation from the coast arc as:
\[
  \delta\xvec(t_f) \;\approx\; \sum_{i=1}^{N} q_i\, \bvec b_i.
\]

\subsection{QUBO formulation}

Let $\bvec d = \xvec_{\mathrm{target}} - \xvec_{\mathrm{coast}}(t_f)$
be the state gap that the burns must close.  The objective balances
trajectory accuracy against fuel usage:
\begin{equation}\label{eq:scheduling-qubo}
  f(\bvec q) \;=\;
  \underbrace{
    \left\|\sum_{i=1}^{N} q_i\, \bvec b_i - \bvec d\right\|_{\bvec W}^2
  }_{\text{miss distance}}
  \;+\;
  \underbrace{
    \alpha \sum_{i=1}^{N} q_i
  }_{\text{fuel penalty}},
\end{equation}
where $\bvec W = \mathrm{diag}(w_1, w_2, w_3, w_4)$ allows
per-component weighting (e.g.\ position only) and $\alpha \ge 0$
penalizes the number of burn segments.

Expanding the quadratic norm and using $q_i^2 = q_i$ for binary
variables:
\begin{align}
  Q_{ij} &= \bvec b_i\TT \bvec W\, \bvec b_j,
    &&\text{quadratic coupling (symmetric)},
    \label{eq:Q-scheduling}\\
  \ell_i &= \alpha - 2\,\bvec b_i\TT \bvec W\, \bvec d,
    &&\text{linear term},
    \label{eq:ell-scheduling}\\
  c &= \bvec d\TT \bvec W\, \bvec d,
    &&\text{constant offset.}
    \label{eq:const-scheduling}
\end{align}

The QUBO energy is then
$f(\bvec q) = \bvec q\TT \Qmat\, \bvec q + \bm\ell\TT \bvec q + c$.

\begin{remark}
Unlike the continuous-lstsq QUBO of Section~5, there is no
quantization error whatsoever: every $q_i$ \emph{is} a physical
binary decision, and the QUBO exactly encodes the scheduling
objective \eqref{eq:scheduling-qubo}.  The only approximation is
the linearization of thrust effects via the STM, which is standard
practice in preliminary mission design.
\end{remark}

\subsection{Solvers}

\paragraph{Brute force.}
For $N \le 20$ ($2^{20} \approx 10^6$ evaluations), the global
optimum is found by exhaustive enumeration.  This provides a ground
truth against which heuristic solvers are validated.

\paragraph{Simulated annealing (SA).}
The QUBO is converted to a \texttt{dimod.BinaryQuadraticModel} and
sampled using \texttt{neal.SimulatedAnnealingSampler}.  SA matches
the brute-force optimum on all problem sizes tested ($N = 8$ to $20$)
with 500--1000 reads.

\paragraph{D-Wave quantum processing unit (QPU).}
The scheduling QUBO maps directly to D-Wave hardware: each decision
step is one physical qubit, and the $Q_{ij}$ couplings are programmed
onto the Pegasus/Zephyr topology.  With $N$ qubits, the problem fits
comfortably within the Advantage2 processor's 5000+ qubits, even for
$N$ in the hundreds.

\paragraph{Mixed-integer linear programming (MILP) baseline.}
A classical, provably optimal baseline is obtained by lifting the
binary-quadratic objective to a pure MILP via the McCormick
linearization
\(
  y_{ij} \le q_i,\;
  y_{ij} \le q_j,\;
  y_{ij} \ge q_i + q_j - 1,\;
  y_{ij} \in \{0,1\},
\)
which exactly enforces $y_{ij} = q_i q_j$.  Substituting
$q_i^2 = q_i$ for binary variables, the QUBO energy becomes
\[
  E(\bvec q) = \sum_{i} (Q_{ii} + l_i)\, q_i
             + \sum_{i<j} 2\, Q_{ij}\, y_{ij} + \mathrm{const},
\]
which is linear in $(\bvec q, \bvec y)$ and solvable by HiGHS through
\texttt{scipy.optimize.milp}.  This provides certified optimality for
small to medium $N$ and a neutral classical scaling reference.

\subsection{Iterative re-linearization}\label{sec:iterative-relinearization}

The QUBO uses a single first-order STM expansion around a reference
trajectory.  For a long arc or a strong thrust the linearization
error accumulates: the schedule that minimizes the linearized
objective need not minimize the true nonlinear miss, and may even
\emph{increase} it.  We address this with a sequential
re-linearization driver, the binary analogue of the continuous-side
Pass-2 loop in Section~\ref{sec:sequential-linearization}.

Let $\bvec q^{(k)}$ be the schedule at iteration $k$ and let
$\bvec x^{(k)}_{\mathrm{nom}}(t)$ be the trajectory obtained by
propagating the true nonlinear CR3BP under the controls implied by
$\bvec q^{(k)}$.  We linearize the final state as
\begin{equation}\label{eq:lin-around-q}
  \bvec x(t_f) \;\approx\; \bvec x^{(k)}_{\mathrm{nom}}(t_f)
  \;+\; \sum_j (q_j - q^{(k)}_j)\, \bvec b^{(k)}_j,
\end{equation}
where $\bvec b^{(k)}_j$ is the impulse-response vector evaluated using
STMs along $\bvec x^{(k)}_{\mathrm{nom}}(t)$.  Substituting into the
miss-distance penalty and absorbing the linear shift into an
\emph{effective} target gap
\[
  \bvec d_{\mathrm{eff}}^{(k)}
  = \big(\bvec x_{\mathrm{target}} - \bvec x^{(k)}_{\mathrm{nom}}(t_f)\big)
  + \sum_j q^{(k)}_j\, \bvec b^{(k)}_j ,
\]
yields a QUBO of the same algebraic form as
\eqref{eq:multi-qubo} with $\bvec d \to \bvec d_{\mathrm{eff}}^{(k)}$
and $\bvec b \to \bvec b^{(k)}$.

The outer loop is then:
\begin{enumerate}[leftmargin=1.5em, itemsep=1pt]
  \item Initialize $\bvec q^{(0)} = \bvec 0$ (coast).
  \item Build the QUBO linearized around $\bvec q^{(k)}$.
  \item Sample a candidate $\hat{\bvec q}$ (any sampler: BF, SA, MILP,
    Leap Hybrid, QPU).
  \item Propagate $\hat{\bvec q}$ under the true nonlinear dynamics
    and measure the actual miss
    $\|\bvec x(t_f) - \bvec x_{\mathrm{target}}\|$.
  \item \emph{Trust-region accept:} if the candidate's true miss
    decreases, accept ($\bvec q^{(k+1)} \leftarrow \hat{\bvec q}$);
    else terminate.
  \item Stop when the schedule is a fixed point or the miss
    improvement falls below tolerance.
\end{enumerate}

Step~5 is the key safeguard.  The linearized model is only locally
accurate; a candidate that the QUBO predicts to be better may in fact
worsen the true miss because the linearization is loose.  Rejecting
non-improving candidates makes the loop monotone in the truth metric
and guarantees that the returned schedule is no worse than coast — a
property the single-pass solve does not have.

\subsection{Multi-direction scheduling}\label{sec:multi-direction}

The single-channel formulation restricts thrust to a single direction
(e.g.\ tangential) at every step.  In practice, a spacecraft may need
to combine thrust components along different axes.  The
\emph{multi-channel} extension introduces $K$ independent thrust
directions, each with its own on/off binary variable per step.

\paragraph{Variable structure.}
Let the thrust channels be labelled $c = 1, \dots, K$ (e.g.\
$K = 2$ for tangential and normal).  At each decision step $i$, the
QUBO contains one variable $q_{c,i} \in \{0,1\}$ per channel.  The
total number of binary variables is $M = K \cdot N$.  Variables are
ordered channel-major:
\[
  \bvec q = (\underbrace{q_{1,1}, \dots, q_{1,N}}_{\text{channel 1}},\;
             \underbrace{q_{2,1}, \dots, q_{2,N}}_{\text{channel 2}},\;
             \dots)
  \;\in\; \{0,1\}^M.
\]

\paragraph{Available directions.}
Four velocity-relative directions are supported:
\begin{itemize}[leftmargin=1.5em, itemsep=1pt]
  \item \textbf{tangential} --- along $\bvec v / \|\bvec v\|$,
  \item \textbf{anti-tangential} --- opposite to $\bvec v$,
  \item \textbf{normal} --- perpendicular to $\bvec v$, rotated
    $90^\circ$ counter-clockwise,
  \item \textbf{anti-normal} --- opposite to normal.
\end{itemize}
These are evaluated at the coast-arc state at each decision point, so
the direction naturally follows the spacecraft's trajectory.

\paragraph{QUBO assembly.}
Each (channel, step) pair contributes an impulse response vector
$\bvec b_{c,i} = \Phi(t_i, t_f)\, \bvec B_{\mathrm{ctrl}}\,
\uvec_{c,i}\, \Delta t$, where $\uvec_{c,i}$ is the thrust
acceleration vector for channel $c$ at step $i$.  The QUBO objective
generalizes naturally:
\begin{equation}\label{eq:multi-qubo}
  f(\bvec q) = \left\|
    \sum_{c=1}^{K} \sum_{i=1}^{N} q_{c,i}\, \bvec b_{c,i} - \bvec d
  \right\|_{\bvec W}^2
  + \alpha \sum_{c,i} q_{c,i}.
\end{equation}
The $M \times M$ QUBO matrix $\Qmat$ has the same structure as
\eqref{eq:Q-scheduling}--\eqref{eq:const-scheduling}, with the
summation indices running over all $(c, i)$ pairs.

\paragraph{Mutual exclusion constraint.}
If the spacecraft has a single engine that cannot fire in two
directions simultaneously, an exclusion penalty is added:
\[
  f_{\mathrm{excl}}(\bvec q) = P \sum_{i=1}^{N} \sum_{c < c'}
  q_{c,i}\, q_{c',i},
\]
where $P$ is a penalty weight chosen large enough to make
simultaneous firing energetically unfavorable.  This adds $P$ to the
off-diagonal QUBO entries coupling channels at the same step, and
preserves the QUBO structure — no auxiliary variables are needed.

\paragraph{Scaling.}
With $K$ channels and $N$ steps, the QUBO has $M = KN$ variables.
For $K = 2$ (tangential $+$ normal), a 50-step problem requires
100~qubits — still well within hardware limits.  The search space
grows to $2^{KN}$, making the combinatorial advantage of quantum
annealing even more relevant than in the single-channel case.


% ==================================================================
\section{Scheduling results}\label{sec:scheduling-results}
% ==================================================================

\subsection{Validation: SA matches brute force}

A cislunar transfer scenario was constructed with initial state
near Earth, a target generated by continuous thrust propagation, and
$N = 15$ decision steps.  The QUBO was solved by both brute force and
simulated annealing (1000 reads).  SA recovered the exact brute-force
optimum in every run, confirming that the QUBO energy landscape is
well-conditioned for annealing-based samplers.

\subsection{Scaling analysis}

The critical motivation for quantum annealing is scaling.
Table~\ref{tab:scheduling-scaling} reports solve times for brute force
and SA as $N$ increases.

\begin{table}[h]
\centering
\begin{tabular}{rrrrl}
  \toprule
  $N$ & $2^N$ & BF (ms) & SA (ms) & SA = BF? \\
  \midrule
   8 &       256 &      3 &    133 & Y \\
  10 &     1\,024 &      8 &    137 & Y \\
  12 &     4\,096 &     30 &    148 & Y \\
  14 &    16\,384 &    116 &    159 & Y \\
  16 &    65\,536 &    435 &    181 & Y \\
  18 &   262\,144 &  1\,920 &    215 & Y \\
  20 & 1\,048\,576 & 10\,100 &    249 & Y \\
  \bottomrule
\end{tabular}
\caption{Brute-force vs.\ simulated-annealing solve times for
  the on/off scheduling QUBO.  SA matches the brute-force global
  optimum at every size tested.  BF time doubles with each additional
  step ($O(2^N)$); SA grows slowly.}
\label{tab:scheduling-scaling}
\end{table}

Key observations:
\begin{itemize}[leftmargin=1.5em, itemsep=2pt]
  \item \textbf{Brute force} scales as $O(2^N)$: doubling time with
    each additional decision step.  At $N=20$ it takes $\sim 10$~s;
    at $N=30$ the extrapolated time exceeds $10^4$~s ($\sim 3$~hours);
    at $N=50$ it is astronomically large ($\sim 10^{15}$ evaluations).
  \item \textbf{Simulated annealing} grows subexponentially in this
    regime and matches the brute-force optimum for all $N \le 20$.
    However, SA is a heuristic: there is no guarantee it will find the
    global minimum for larger $N$ with strong couplings.
  \item \textbf{D-Wave QPU} anneal time is $\sim 20$~\si{\micro\second}
    regardless of $N$, as long as the QUBO fits on the hardware graph.
    A 50-step scheduling problem requires 50~qubits, well within the
    Advantage2 processor's $>5000$-qubit capacity.
\end{itemize}

\subsection{Where quantum advantage emerges}

The regime $N \gtrsim 30$ is where quantum annealing becomes
compelling:
\begin{itemize}[leftmargin=1.5em, itemsep=2pt]
  \item Classical brute force is intractable ($2^{30} > 10^9$).
  \item SA may find good solutions but lacks optimality guarantees.
  \item The QPU solves the \emph{exact same} QUBO in
    $\sim 20$~\si{\micro\second}, with the annealing physics
    naturally exploring the combinatorial landscape.
\end{itemize}
This is a qualitative difference from the continuous-lstsq QUBO, where
the QPU was solving a worse approximation of a problem that classical
methods handle exactly.  Here, the QPU solves the \emph{exact}
problem, and it does so faster.

\subsection{Multi-direction results}

Table~\ref{tab:multi-direction} compares single-channel (tangential
only) and multi-channel (tangential~$+$~normal) scheduling on a
lateral-correction scenario with $N = 10$ decision steps and
position-only targeting ($W = \mathrm{diag}(1,1,0,0)$).
The spacecraft departs from a high Earth orbit ($r \approx 200\,000$~km)
and the target is generated by propagating with pure normal thrust
($a = 0.02$, $t \in [0, 1]$), creating a gap that tangential burns
alone cannot close effectively.

\begin{table}[h]
\centering
\begin{tabular}{lrrrl}
  \toprule
  Mode & Qubits & Burns & Position miss & Improvement \\
  \midrule
  Tangential only         & 10 &  1 & $1.02 \times 10^{-2}$ & (baseline) \\
  Tangential $+$ Normal   & 20 &  9 & $2.61 \times 10^{-5}$ & $-99.7\%$ \\
  Exclusive (tan $+$ nor) & 20 &  9 & $8.18 \times 10^{-5}$ & $-99.2\%$ \\
  \bottomrule
\end{tabular}
\caption{Comparison of single- and multi-channel scheduling on a
  lateral-correction scenario.  The multi-channel QUBO combines
  tangential and normal burns to reduce the position miss by nearly
  three orders of magnitude.  The exclusive variant, which enforces at
  most one channel per step via a quadratic penalty, still achieves
  $99.2\%$ improvement by alternating between thrust directions.}
\label{tab:multi-direction}
\end{table}

The multi-channel advantage emerges when the target requires
corrections outside the velocity direction.  The scenario uses
position-only targeting, which is physically realistic for waypoint
problems where velocity can be corrected by a subsequent maneuver.
The key insight is that the STM couples all state components over
long arcs, washing out the distinction between thrust directions.
A shorter transfer arc ($t_f = 1$ vs.\ $t_f = 3$) preserves the
directional separation and allows normal burns to retain their
lateral character at the final time.

The multi-channel formulation scales naturally:
With $K = 2$ channels and $N = 50$ steps, the QUBO has
$M = 100$~binary variables — a $2^{100}$ search space that is utterly
intractable by brute force but fits on 100~physical qubits.

\subsection{Linearization-gap and iterative correction}\label{sec:lin-gap}

The single-pass QUBO assumes the first-order STM expansion is
accurate over the whole time interval.  In a long-arc cislunar
correction at $T = 4$ nondim ($N = 15$,
$v_\mathrm{inj}/v_\mathrm{circ} = 1.187$) the assumption fails:
the QUBO's optimum schedule, when propagated under the true
nonlinear CR3BP, lands $\sim 143\,000$~km from the target, almost
identical to coast and far from the linearised
prediction.  The iterative re-linearisation driver
(Section~\ref{sec:iterative-relinearization}) re-builds the QUBO
around the current schedule at each outer iteration and applies
a trust-region accept rule; on this scenario it converges in four
outer iterations to a schedule whose true miss is $\sim 4\,800$~km,
a $30 \times$ reduction at the cost of one extra burn.
Figure~\ref{fig:lin-gap} compares the two trajectories side by
side.

\begin{figure}[htbp]
  \centering
  \includegraphics[width=0.98\textwidth]{../scripts/iterative_vs_singlepass_trajectory.png}
  \caption{Single-pass QUBO (left panel) versus iterative
    re-linearisation (right panel) on a long-arc cislunar
    correction at $T = 4$ nondim. The blue star is the target
    final state; the cross is each schedule's achieved final
    state. The single-pass scheduler picks two burns (red
    segments) whose linearised effect is correct but whose
    nonlinear effect leaves the spacecraft on a markedly
    different arc. The iterative scheduler converges in four
    outer iterations to a three-burn schedule that lands within
    $4\,800$~km of the target.}
  \label{fig:lin-gap}
\end{figure}

For arcs longer than $T \approx 5$ nondim the linearisation is so
loose that the first candidate is non-improving and the iterative
loop terminates at coast --- it \emph{abstains} rather than
acting on bad information.  The single-pass solve, lacking the
trust-region rule, would in this regime recommend a schedule that
is worse than doing nothing.

\subsection{Second application: lunar-orbit station-keeping}\label{sec:moon-orbit}

The same iterative scheduler applies to mission phases beyond
injection correction.  Figure~\ref{fig:moon-orbit} shows a
$5\,000$-km circular lunar orbit perturbed by a $0.5\%$
tangential velocity error at $t_0$.  Without correction the
spacecraft drifts $\sim 1\,300$~km from the nominal endpoint
after two revolutions; the iterative QUBO selects a single
anti-tangential channel and uses $4.4$~m/s of $\Delta v$ to
reduce the position miss to $\sim 530$~km.

\begin{figure}[htbp]
  \centering
  \includegraphics[width=0.98\textwidth]{../scripts/moon_orbit_stationkeeping.png}
  \caption{Lunar-orbit station-keeping with the iterative QUBO
    at the same spacecraft class as the injection-correction
    scenario (100-kg microsatellite, $5.4$-mN ion thruster).
    Left: synodic frame zoomed near the Moon over two
    revolutions ($\sim 28$~hours).  Right: Moon-relative frame
    in km.  The QUBO uses position-only target weighting
    ($W = \mathrm{diag}(1,1,0,0)$), so a phase-shift along the
    orbit affects the velocity vector but not the position
    objective --- the right metric for a station-keeping
    mission that cares about position, not phase.}
  \label{fig:moon-orbit}
\end{figure}

\subsection{Classical MILP baseline}

The MILP lift described in Section~\ref{sec:thrust-scheduling}
provides a certified-optimal classical reference at small to medium
$N$.  Solved through HiGHS via \texttt{scipy.optimize.milp}, the
lift produces an exact optimum at $N = 20$ in $\sim 12$~s with
optimality certificate; at $N = 30$ the same solver cannot certify
optimality within a 120-second budget but returns a feasible
incumbent indistinguishable from SA.  The McCormick lift creates
$N(N{-}1)/2$ auxiliary binaries and a dense LP relaxation, so the
classical-MILP scaling crosses the simulated-annealing curve in the
$N \approx 25$--$30$ region.  This is the operationally relevant
crossover: where SA dominates a state-of-the-art classical
exact solver, the case for moving to D-Wave hardware becomes
quantitatively defensible.


% ==================================================================
\section{Multi-arc capture and end-to-end mission}\label{sec:multi-arc}
% ==================================================================

The iterative re-linearisation of
Section~\ref{sec:iterative-relinearization} closes the
linearisation gap of a \emph{single} arc. To complete the mission
(spiral out from GEO, correct the cislunar transfer, and capture
into a bound lunar orbit) we extend the scheduler in two ways:
a sliding-window structure that chains short single-arc QUBOs
across many lunar revolutions, and a composition-of-phases
structure that switches the action set, target weighting and
thrust regime as the spacecraft moves between Earth-, cislunar-
and Moon-dominated regimes.

% ------------------------------------------------------------------
\subsection{Sliding-window QUBO for lunar capture (Phase 3)}\label{sec:sliding-window}
% ------------------------------------------------------------------

A representative case the single-arc QUBO cannot handle is lunar
orbit insertion (LOI) from a hyperbolic flyby: the spacecraft
passes the Moon once, must shed several tens of metres per second
of excess velocity to bind its orbit, and then circularise over
many subsequent revolutions. A single-arc QUBO covering this entire
process would have a linearisation gap so large that no schedule
survives the trust-region rule of
Section~\ref{sec:iterative-relinearization}. The implementation
lives in \texttt{qalunar/qubo/lunar\_capture.py} and is exposed as
\texttt{solve\_lunar\_capture\_sliding\_window}.

\paragraph{Sliding-window structure.} We split the mission into $W$
short arcs of duration $\Delta T_w$, each of which is small enough
that the iterative re-linearisation converges. Between windows the
spacecraft state is propagated under the true nonlinear CR3BP and
becomes the initial state of the next window. The total schedule
is the concatenation of per-window schedules.

\paragraph{Window-length cap.} The natural choice for $\Delta T_w$
is a fraction of the local Moon-relative orbital period
$T_\mathrm{loc} = 2\pi\sqrt{r_w^3/\mu}$, with $r_w$ the
Moon-relative radius at the start of window~$w$. This adapts the
window to the local dynamics. However, $T_\mathrm{loc}$ diverges
as $r_w$ grows (loose orbits or escapes), and an unguarded
fractional-period window would become arbitrarily long once the
orbit is large -- by which point the linearisation has long since
collapsed. We therefore cap $\Delta T_w$ at a fixed
$\Delta T_\mathrm{max}$ so the linearisation regime stays valid
even when the orbit is loose.

\paragraph{Action-radius gating.} A second safeguard restricts QUBO
scheduling to windows in which the spacecraft is within a threshold
$r_\mathrm{act}$ of the Moon. Outside this radius, dynamics are
dominated by Earth and the Moon-relative linear target is
meaningless: every candidate fails the trust-region rule and the
loop produces no progress while burning through windows. With
gating, those windows become unforced drifts -- the spacecraft is
propagated freely under the CR3BP until it returns inside
$r_\mathrm{act}$, at which point QUBO scheduling resumes. Choosing
$r_\mathrm{act}$ on the order of a few perilune radii confines the
QUBO action to the region where anti-tangential thrust efficiently
dissipates orbital energy.

\paragraph{Per-window target.} For each window, we set the QUBO
target to a state whose Moon-relative inertial velocity has
magnitude $v_\mathrm{circ}(r) = \sqrt{\mu/r}$ in the direction of
the spacecraft's current Moon-relative velocity, evaluated at the
unforced end-of-window position $\xvec_\mathrm{eow}$:
\begin{equation}\label{eq:sw-target}
  \xvec_\mathrm{target}^{(w)}
  = \big[\,
      \bvec{r}_\mathrm{eow},\;
      v_\mathrm{circ}(\|\bvec{r}_\mathrm{eow} - \bvec{r}_\mathrm{Moon}\|)\,
      \hat{\bvec{v}}_\mathrm{Moon\text{-}rel}^{\,\mathrm{eow}}
    \,\big]
\end{equation}
expressed in the synodic frame after the standard $\bvec{\omega}\times\bvec{r}$
correction. Position is left unweighted ($W_\mathrm{pos} = 0$) and only
velocity components are weighted in the QUBO objective: the target
says \emph{stay where the unforced trajectory takes you, but slow to
the local circular speed}. With a single anti-tangential channel,
this target elicits firings precisely where the spacecraft is moving
faster than circular -- i.e.\ near perilune of an elliptical pass --
which is the energy-dissipation regime for lunar capture and
circularisation. At apolune the spacecraft is already slower than
$v_\mathrm{circ}$, anti-tangential thrust cannot speed it up, and
the QUBO correctly stays off.

\paragraph{Termination and diagnostics.} The outer loop terminates
after $W$ windows have been processed. We monitor the two-body
energy
$E_w = \tfrac{1}{2}\|\bvec{v}_w^\mathrm{rel}\|^2 -
\mu/\|\bvec{r}_w^\mathrm{rel}\|$
and the Moon-relative apolune radius $r_\mathrm{apo}$ across
windows. Captured states have $E_w < 0$ and finite
$r_\mathrm{apo}$; loose-orbit windows have $r_\mathrm{apo}$ in the
tens of thousands of kilometres; escaped or pre-capture states
have $E_w \ge 0$ and $r_\mathrm{apo} = \infty$.

\paragraph{Per-window guarantees.} Each window is itself a call to
the iterative driver of
Section~\ref{sec:iterative-relinearization} with a short
$\Delta T_w$, so the trust-region monotonicity property applies
window-by-window: the schedule produced for window~$w$ has true
nonlinear miss no worse than coast \emph{within that window}. The
composition does not admit a global monotonicity guarantee in
$E_w$ because the unforced drift between windows can both raise or
lower the energy under CR3BP perturbations, but in practice
anti-tangential firings near perilune dominate and the energy
decreases over the course of the schedule (see
Section~\ref{sec:end-to-end-manual}).

\paragraph{Sensitivity sweep.} The three principal parameters
($\Delta T_\mathrm{max}$, $r_\mathrm{act}$, brake factor on
$v_\mathrm{circ}$) trade off as expected: larger
$\Delta T_\mathrm{max}$ allows more burn slots per window but
worsens the in-window linearisation; larger $r_\mathrm{act}$
allows scheduling further from the Moon at the cost of less
effective burns; a brake factor below $1$ targets a sub-circular
speed and produces lower-apolune captures at higher $\Delta v$
cost. The sweep is reproduced by
\texttt{scripts/run\_sliding\_window\_sweep.py} and reported in
the companion paper.

% ------------------------------------------------------------------
\subsection{End-to-end mission: composing Phases 1, 2 and 3}\label{sec:end-to-end-manual}
% ------------------------------------------------------------------

The end-to-end demo
(\texttt{scripts/run\_full\_mission\_demo.py} for numbers,
\texttt{scripts/run\_mission\_figure.py} and
\texttt{run\_continuity\_figure.py} for figures) flies a single
\SI{100}{\kg} microsatellite with a single throttleable ion/Hall
thruster (\SI{350}{\milli\newton} main, \SI{5.4}{\milli\newton}
cruise) from GEO to a bound near-circular lunar orbit, with no
impulsive burns. The three phases share the same spacecraft and
the same engine; only the throttling state, the action channel
and the target weighting change between phases.

\paragraph{Phase 1 -- Earth ascent + injection boost.} The
analytical Edelbaum spiral of
Section~\ref{sec:thrust-scheduling} (implemented in
\texttt{qalunar/reference/edelbaum.py}) lifts the spacecraft from
GEO ($r \approx \SI{42164}{\km}$) to a circular parking orbit at
$r = \SI{200000}{\km}$ at \emph{main} thrust. A continuous
all-on tangential arc on the same thruster then adds the
$\SI{264}{\meter\per\second}$ needed to inject onto a Moon-bound
trajectory; this corresponds to \SI{20.9}{\hour} of thrust and
replaces what would otherwise be an impulsive chemical burn.
Phase~1 totals $\Delta v \approx \SI{1924}{\meter\per\second}$
over \SI{6.4}{\day}.

\paragraph{Phase 2 -- cislunar correction.} The thruster throttles
to cruise so each binary slot delivers the
$\sim \SI{4.4}{\meter\per\second}$ needed to discriminate
trajectory corrections at the metres-per-second level. This phase
is exactly the long-arc scenario validated in
Section~\ref{sec:lin-gap}: $T = 4$ nondim, $N = 15$, iterative
re-linearisation, $\Delta v \approx \SI{16}{\meter\per\second}$
in three burns over \SI{17.4}{\day}, leaving the spacecraft on a
hyperbolic flyby with perilune at $\sim$~\SI{4500}{\km} altitude.

\paragraph{Phase 3 -- sliding-window lunar capture.} The thruster
returns to main power and the sliding-window scheduler of
Section~\ref{sec:sliding-window} runs from the raw Phase-2
perilune state -- no impulsive pre-capture burn. With
$\Delta T_\mathrm{max} = 0.01$ nondim ($\sim$~\SI{1}{\hour}),
$N = 12$ decision steps per window, an action radius of
$3\,r_\mathrm{perilune}$, and a single anti-tangential channel,
the loop fires during the close approach to the Moon. After
$\sim$~\num{200} windows the energy has dropped from $E = +0.087$
(hyperbolic) to $E = -0.38$ (bound) and the apolune has collapsed
from infinity to $r_\mathrm{apo} \approx \SI{6212}{\km}$, a
near-circular bound orbit at the perilune radius. Phase~3 expends
$\Delta v \approx \SI{561}{\meter\per\second}$ in \SI{8}{\day}.

\paragraph{Mission totals.} The complete mission delivers
$\Delta v_\mathrm{total} \approx \SI{2502}{\meter\per\second}$ in
\SI{31.8}{\day} of flight, with every metre per second produced by
binary on/off commands of a single throttleable thruster
(Table~\ref{tab:end-to-end-manual}). Phase~1 dominates the
$\Delta v$ budget at \(77\%\) of the total, Phase~3 is the
second-largest cost at \(22\%\), and Phase~2 contributes
\(0.6\%\), essentially free in comparison; yet without the
Phase-2 correction the \num{22000}-km miss precludes any usable
perilune passage and the mission cannot continue.

\begin{table}[htbp]
\centering
\caption{All-binary Earth--Moon mission breakdown. Single
\SI{100}{\kg} microsatellite with a throttleable ion/Hall thruster
(\SI{350}{\milli\newton} main / \SI{5.4}{\milli\newton} cruise).
No impulsive burns.}
\label{tab:end-to-end-manual}
\begin{tabular}{lrrl}
\toprule
Phase & $\Delta v$ (m/s) & TOF (days) & Mode \\
\midrule
1: Edelbaum spiral + injection boost
  & \num{1924} & \num{6.4} & main, tangential \\
2: cislunar correction (QUBO)
  & \num{16}   & \num{17.4} & cruise, tangential \\
3: sliding-window QUBO capture
  & \num{561}  & \num{8.0}  & main, anti-tangential \\
\midrule
Total
  & \num{2502} & \num{31.8} & all binary \\
\bottomrule
\end{tabular}
\end{table}

\paragraph{Why throttling is necessary.} Bumping every phase to
main thrust breaks Phase~2: at \SI{350}{\milli\newton} each binary
slot delivers $\sim \SI{350}{\meter\per\second}$, vastly more
than the metres-per-second resolution needed to discriminate good
trajectory corrections. The QUBO either fires zero slots
(under-correction) or a single slot (overshooting by an order of
magnitude). Throttling to \SI{5.4}{\milli\newton} restores the
fine-grained regime that Phase~2 requires.

\paragraph{Why the action-radius gate is necessary.} Without
gating, the sliding-window solver fires the QUBO during windows in
which the spacecraft is close to L1 and dominated by Earth's
gravity. The Moon-relative circular target has no physical meaning
in that regime; the trust-region rule rejects every candidate, but
more dangerously, when a candidate is accidentally accepted
(because the linearised cost happens to look better) the true
CR3BP propagation perturbs the orbit further. With the gate, the
spacecraft drifts naturally between perilune passages and QUBO
firings concentrate where they are physically effective.

% ------------------------------------------------------------------
\subsection{Connection to primer-vector theory}\label{sec:primer-vector-manual}
% ------------------------------------------------------------------

The on/off scheduling QUBO never computes a Pontryagin costate ---
the decision variables are the binary controls themselves, not the
adjoint. Once a schedule has converged under the iterative driver,
however, a costate can be reconstructed \emph{a posteriori} along
the resulting nonlinear trajectory by integrating
$\dot{\bvec{\lambda}} = -A(\xvec(t))^\top \bvec{\lambda}$
backward in time from the terminal transversality condition
$\bvec{\lambda}(t_f) = 2\bvec{W}(\xvec_f - \xvec_\mathrm{target})$.
The tangential primer-vector analogue is the scalar switching
function
\begin{equation}\label{eq:switching-manual}
  S(t) = -\bvec{\lambda}_v(t) \cdot \hat{\bvec{v}}(t),
\end{equation}
where $\bvec{\lambda}_v$ is the velocity sub-vector of the costate
and $\hat{\bvec{v}}$ the unit velocity vector. Pontryagin's minimum
principle predicts that the optimal tangential bang-bang control
fires the engine when $S > 0$ and coasts otherwise. Comparing this
continuous-time prediction with the QUBO's discrete schedule
provides a sanity check on the bridge between the binary
formulation and indirect primer-vector
methods~\cite{Taylor2026}. The implementation is in
\texttt{qalunar/reference/primer\_vector.py} and the figure is
produced by \texttt{scripts/run\_primer\_vector\_figure.py}.

Applied to the converged iterative schedule of
Section~\ref{sec:lin-gap} ($T = 4$ nondim, $N = 15$, three ON
slots out of fifteen), the reconstruction yields a switching
function whose support is a contiguous interval near the second
half of the arc. Three observations stand out:
\begin{enumerate}[leftmargin=1.5em]
  \item Every QUBO ON-slot has its slot centre inside the
    PMP-predicted ON region ($S(t_i) > 0$, $3/3$); the converse
    pathology --- a QUBO burn placed where PMP predicts coast ---
    does not occur.
  \item The QUBO does not saturate the $S > 0$ interval as a
    continuous PMP would in the unconstrained case: nine of the
    fifteen slot centres lie in the $S > 0$ region but only three
    are activated. With no fuel penalty, continuous PMP would fire
    the engine across the entire $S > 0$ interval; the binary
    QUBO instead selects a minimal subset of those slots
    sufficient to land the target within tolerance, because the
    squared-norm objective is minimised by the smallest set of
    impulse-response vectors $\bvec{b}_i$ whose sum cancels the
    gap $\bvec{d}_\mathrm{eff}$, not by activating every slot
    whose individual contribution would align with
    $-\bvec{d}_\mathrm{eff}$.
  \item The costate exhibits strong backward amplification near
    $t_0$, a direct manifestation of CR3BP sensitivity to the
    initial state: small perturbations early in the arc compound
    through the chaotic flow and produce large terminal-state
    changes, which is exactly the regime where the linearisation
    gap is widest --- and exactly why the iterative loop is
    required to make the QUBO faithful in the first place.
\end{enumerate}

This reconstruction is a self-consistency check:
$\bvec{\lambda}$ is built from the QUBO's own final state and
propagated through the QUBO's own trajectory. A fully external
comparison would solve the same problem via indirect methods (the
route taken in~\cite{Taylor2026} via homotopy continuation in
$\rho$ and $u_\mathrm{max}$) and overlay the two switching
functions. We leave that comparison to future work, as the
in-house direct-collocation NLP referenced earlier is
energy-optimal rather than fuel-optimal and therefore lacks a
directly comparable primer vector.

% ------------------------------------------------------------------
\subsection{Comparison with state-of-the-art Earth--Moon transfers}\label{sec:purdue-comparison}
% ------------------------------------------------------------------

Recent work by Taylor et~al.~\cite{Taylor2026} from the Purdue
Orbital student team designs Earth-to-NRHO and Earth-to-LLO
transfers via Pontryagin's minimum principle and primer-vector
theory, with hyperbolic-tangent smoothing and homotopy continuation
in $\rho$ and $u_\mathrm{max}$ to converge on the impulsive limit.
Their reported totals are \SI{3.93}{\km\per\second} in
\SI{5.2}{\day} for LEO\,$\to$\,NRHO and \SI{3.91}{\km\per\second}
in \SI{4.7}{\day} for LEO\,$\to$\,LLO. Table~\ref{tab:purdue-manual}
places those numbers next to the qalunar end-to-end binary-QUBO
result.

\begin{table}[htbp]
\centering
\caption{End-to-end Earth--Moon transfer comparison. The
\emph{method} column lists the principal optimisation framework;
totals are not directly comparable due to differences in starting
orbit, propulsion system and final orbit (see text).}
\label{tab:purdue-manual}
\begin{tabular}{llrrl}
\toprule
Source & Transfer & $\Delta v$ (km/s) & TOF (d) & Method \\
\midrule
This work
  & GEO\,$\to$\,bound lunar orbit
  & \num{2.50} & \num{31.8}
  & binary QUBO, low-thrust ion/Hall \\
\cite{Taylor2026}
  & LEO\,$\to$\,9{:}2 NRHO
  & \num{3.93} & \num{5.2}
  & PMP + primer vector, impulsive chemical \\
\cite{Taylor2026}
  & LEO\,$\to$\,LLO (\SI{100}{\km})
  & \num{3.91} & \num{4.7}
  & PMP + primer vector, impulsive chemical \\
\bottomrule
\end{tabular}
\end{table}

The comparison is not apples-to-apples in three respects:
(i)~the starting orbit differs (GEO vs.\ LEO --- a
$\sim\SI{3.9}{\km\per\second}$ Hohmann gap that favours qalunar);
(ii)~the propulsion system differs (low-thrust electric vs.\
impulsive chemical --- low thrust buys $\Delta v$ efficiency at
the cost of TOF); and (iii)~the final orbit differs (qalunar
reaches a near-circular bound orbit at the perilune radius rather
than \SI{100}{\km} LLO or the 9{:}2 NRHO). Subject to those
caveats, the qualitative conclusion is that the binary-QUBO
formulation produces a $\Delta v$ budget consistent with the
indirect-method state of the art at the cost of an order of
magnitude longer TOF --- the expected low-thrust trade-off. The
methodological novelty is not a $\Delta v$ improvement but the
demonstration that an entire Earth--Moon mission can be flown by
binary on/off commands of a single throttleable thruster, with
the QUBO scheduler and its sliding-window extension as the
discrete-side analogue of the indirect primer-vector switching
function~\cite{Taylor2026}.


% ==================================================================
\section{Discussion: why the project changed direction}\label{sec:discussion-pivot}
% ==================================================================

The original project plan, following Carbone et~al.~\cite{Carbone2025},
aimed to extend the QUBO-encoded continuous least-squares approach
from the two-body HCW setting to the nonlinear CR3BP.
The investigation described in
Section~\ref{sec:convergence-investigation} revealed a fundamental
mismatch between the tool (binary optimization) and the problem
(continuous optimization):

\begin{enumerate}[leftmargin=1.5em]
  \item \textbf{Classical methods solve the inner problem exactly.}
    Least squares finds the continuous minimum in $\sim 5$~ms.
    The QUBO \emph{approximates} the same answer with quantization
    error, in far more time.

  \item \textbf{Quantization error is a hard floor.}
    With $p$~bits per variable, the best the QUBO can achieve
    differs from the true lstsq solution by $O(2^{-p})$ per
    component.  This error propagates through the sequential
    linearization and prevents convergence below a
    quantization-dominated residual plateau.

  \item \textbf{There is no asymptotic advantage.}
    Increasing $p$ shrinks the error but linearly increases the
    number of qubits.  The QUBO is always solving a harder version
    of a problem that classical methods solve trivially.
\end{enumerate}

The pivot to on/off thrust scheduling resolves all three issues.
The decision variables $q_i \in \{0,1\}$ are \emph{inherently binary}:
there is no encoding, no quantization, no approximation beyond the
standard STM linearization.  The $2^N$ search space is genuinely hard
for classical exhaustive methods, and quantum annealing is designed
precisely for this class of combinatorial optimization.

The previous work (TFC+ELM transcription, lstsq-to-QUBO pipeline,
sequential linearization) remains valuable as a reference
implementation and a documented negative result: it shows where QUBO
encoding of continuous problems breaks down and why naturally binary
formulations are the scientifically sound path forward.


% ==================================================================
\section{Glossary of symbols}
% ==================================================================

\begin{center}
\begin{tabular}{llp{7cm}}
  \toprule
  Symbol & Meaning & Notes \\
  \midrule
  $\mu$ & CR3BP mass parameter & $m_2/(m_1+m_2)$; Earth--Moon $\approx 0.01215$ \\
  $r_1, r_2$ & Distances to primaries & $r_1$ from Earth, $r_2$ from Moon \\
  $\Omega$ & Pseudo-potential & Centrifugal $+$ gravitational + bookkeeping \\
  $\Omega_x, \Omega_y$ & Gradient of $\Omega$ & Drives position dynamics \\
  $\Omegaxx, \Omegaxy, \Omegayy$ & Hessian of $\Omega$ & Time-varying along nominal \\
  $C_J$ & Jacobi integral & $2\Omega - \|\bvec{v}\|^2$ \\
  $(L_1, \dots, L_5)$ & Lagrange points & $L_4, L_5$ stable for Earth--Moon \\
  \midrule
  $\bvec{r}, \bvec{v}$ & State position and velocity & $\bvec{r}=(x,y)$, $\bvec{v}=(v_x,v_y)$ \\
  $\uvec$ & Thrust acceleration control & $(u_x, u_y)$ \\
  $\lamr, \lamv$ & Costates & $\lamr$ conjugate to $\bvec r$, $\lamv$ to $\bvec v$ \\
  $\mathcal H$ & Control Hamiltonian & $u^\star = -\lamv$ \\
  $T$ & Time of flight & Fixed \\
  $(\bvec r_0, \bvec v_0)$ & Initial BC & LEO state \\
  $(\bvec r_f, \bvec v_f)$ & Final BC & LLO state \\
  \midrule
  $z \in [-1, 1]$ & Normalized time & $z = c(t - t_0) - 1$, $c = 2/T$ \\
  $\omega_1, \dots, \omega_4$ & Hermite switching functions & TFC BC-annihilators \\
  $\Hmat$ & ELM basis matrix & $n \times L$, $H_{ij} = \tanh(w_j z_i + b_j)$ \\
  $\Fmat$ & Constrained expression matrix & ELM minus switching-projections \\
  $n$ & Training points per coordinate & Typical $40$ \\
  $L$ & ELM neurons per coordinate & Typical $20$--$80$ \\
  $\bm\xi$ & ELM output weights & Free parameters \\
  \midrule
  $\Amat, \Bmat$ & Linear residual system & Shape $6n \times 6L$ and $6n$ \\
  $\lambda_{\mathrm{tik}}$ & Tikhonov strength & $\varepsilon_{\mathrm{tik}} \sigma_{\max}(\Amat)^2$ \\
  $r$ & Effective rank after SVD truncation & Much smaller than $6L$ in practice \\
  $\Vr$ & Right basis of truncated SVD & $6L \times r$ \\
  $\bm\alpha$ & Reduced coordinates & $\bm\xi = \Vr \bm\alpha$ \\
  $p$ & Bits per reduced variable & Typical $8$ \\
  $\bvec q \in \{0,1\}^{rp}$ & QUBO bitstring & Physical qubits (lstsq encoding) \\
  $\Qmat, \bvec\ell, c$ & QUBO triple & $E(\bvec q)=\bvec q^\top \Qmat \bvec q + \bvec\ell^\top \bvec q + c$ \\
  \midrule
  $q_i \in \{0,1\}$ & Thrust on/off decision & $q_i=1$: burn at step $i$ \\
  $N$ & Number of decision steps & $= $ number of qubits (scheduling) \\
  $\Phi(t_i, t_f)$ & State transition matrix & Maps perturbation at $t_i$ to $t_f$ \\
  $\bvec B_{\mathrm{ctrl}}$ & Control influence matrix & $(4\times 2)$, maps thrust to state derivative \\
  $\bvec b_i$ & Impulse response vector & $\Phi(t_i,t_f)\bvec B_{\mathrm{ctrl}} \uvec_i \Delta t$ \\
  $\bvec d$ & Target state gap & $\xvec_{\mathrm{target}} - \xvec_{\mathrm{coast}}(t_f)$ \\
  $\alpha$ & Fuel penalty weight & Trades miss distance vs.\ number of burns \\
  $\bvec W$ & Target weighting matrix & $\mathrm{diag}(w_1,w_2,w_3,w_4)$; default $\bvec I$ \\
  \bottomrule
\end{tabular}
\end{center}


% ==================================================================
\section*{Bibliography}
% ==================================================================
\begin{thebibliography}{9}

  \bibitem{Carbone2025}
  De Grossi, F., Carbone, A., Spiller, D., Ottaviani, D.,
  Mengoni, R., Circi, C.,
  \emph{Transcription and optimization of an interplanetary trajectory through quantum annealing},
  Astrodynamics \textbf{9}:195--215, 2025.

  \bibitem{Szebehely1967}
  Szebehely, V., \emph{Theory of Orbits: The Restricted Problem of
   Three Bodies}, Academic Press, 1967.

  \bibitem{KLMR2011}
  Koon, W., Lo, M., Marsden, J., Ross, S., \emph{Dynamical Systems,
   the Three-Body Problem and Space Mission Design}, 2011.

  \bibitem{Mortari2017}
  Mortari, D., \emph{The Theory of Connections: Connecting Points},
  Mathematics \textbf{5}(4):57, 2017.

  \bibitem{Huang2006}
  Huang, G.-B., Zhu, Q.-Y., Siew, C.-K.,
  \emph{Extreme learning machine: Theory and applications},
  Neurocomputing \textbf{70}:489--501, 2006.

  \bibitem{Kelly2017}
  Kelly, M.,
  \emph{An Introduction to Trajectory Optimization: How to Do Your
   Own Direct Collocation},
  SIAM Review \textbf{59}(4):849--904, 2017.

  \bibitem{Betts2010}
  Betts, J.~T.,
  \emph{Practical Methods for Optimal Control and Estimation Using
   Nonlinear Programming}, 2nd~ed., SIAM, 2010.

  \bibitem{Conn2000}
  Conn, A.~R., Gould, N.~I.~M., Toint, P.~L.,
  \emph{Trust-Region Methods},
  MPS-SIAM Series on Optimization, SIAM, 2000.

  \bibitem{Taylor2026}
  Taylor, T.~E., Madden, K.~F., Bohne, S.~L., Vetter, M.~J.,
  \emph{An Earth--Moon Trajectory and Mission Design},
  Purdue Orbital, Purdue University, West Lafayette, IN, 2026.

\end{thebibliography}

\end{document}
